\documentclass[a4paper,12pt]{article}
\usepackage{docsTemplate}
\usepackage{usecase}
\usepackage{float}

\setTitle{Analisi dei Requisiti}
\setAuthors{Filippo Compagno \\ & Sara Gioia Fichera \\ & Andrea Masiero \\ & Mattia Oliva Medin \\ & Francesco Balestro \\ & Nenad Radulovic \\ & Bianca Zaghetto}
\setVerificators{Andrea Masiero \\ & Francesco Balestro \\ & Bianca Zaghetto \\ & Sara Gioia Fichera \\ & Nenad Radulovic \\ & Filippo Compagno \\ & Mattia Oliva Medin}
\setApprovation{}
\setVersion{0.15}
\setType{Esterno}
\setDestination{Tullio Vardanega \\ & Riccardo Cardin \\ & Sanmarco Informatica \\[-0.1em] & Team Three Technologies}

% -- Obbligatori --
\newcommand{\loaddip}{01}
\newcommand{\remloaddip}{02}
\newcommand{\errloaddip}{03}
\newcommand{\removedip}{04}
\newcommand{\viewdip}{05}
\newcommand{\viewdetdoc}{06}
\newcommand{\viewconsinfo}{07}

% IdDoc
\newcommand{\mdparamImpCrit}{08} % Impronta crittografica - Hash - Algoritmo
\newcommand{\mdparamIdDoc}{09} % Identificativo documento
% ModalitàDiFormazione
\newcommand{\mdparamModFor}{10} % Modalità di formazione
% TipologiaDocumentale
\newcommand{\mdparamTipDoc}{11} % Tipologia documentale
% DatiDiRegistrazione
\newcommand{\mdparamTipFlu}{12} % Tipologia di flusso
\newcommand{\mdparamTipReg}{13} % Tipo registro
\newcommand{\mdparamDatReg}{14} % Data registrazione
\newcommand{\mdparamOraReg}{15} % Ora registrazione
\newcommand{\mdparamNumDoc}{16} % Numero documento
\newcommand{\mdparamCodReg}{17} % Codice registro
\newcommand{\mdparamRuolo}{18} % Ruolo
\newcommand{\mdparamTipSog}{19} % Tipo soggetto
\newcommand{\mdparamCog}{20} % Cognome (PF, AS)
\newcommand{\mdparamNom}{21} % Nome (PF, AS)
\newcommand{\mdparamCF}{22} % Codice fiscale (PF, PG, AS)
\newcommand{\mdparamIndDig}{23} % Indirizzo digitale di riferimento (PF, PG, PAI, PAE, AS)
\newcommand{\mdparamDenOrg}{24} % Denominazione Organizzazione (PG, AS)
\newcommand{\mdparamCFPI}{25} % Partita IVA (PG)
\newcommand{\mdparamDenUff}{26} % Denominazione ufficio (PG, PAE, AS)
\newcommand{\mdparamDenAmm}{27} % Denominazione amministrazione (PAI, PAE)
\newcommand{\mdparamIPA}{28} % Codice IPA (PAI)
\newcommand{\mdparamDenAmmAOO}{29} % Denominazione amministrazione AOO (PAI)
\newcommand{\mdparamIPAOOO}{30} % Codice IPA AOO (PAI)
\newcommand{\mdparamDenAmmUOR}{31} % Denominazione amministrazione UOR (PAI)
\newcommand{\mdparamIPAUOR}{32} % Codice IPA UOR (PAI)
\newcommand{\mdparamDenSis}{33} % Denominazione Sistema (SW)
% ChiaveDescrittiva
\newcommand{\mdparamOgg}{34} % Oggetto
\newcommand{\mdparamParChi}{35} % Parola chiave
% Allegati
\newcommand{\mdparamNumAll}{36} % Numero allegati
\newcommand{\mdparamImpCritAll}{37} % Impronta crittografica allegato
\newcommand{\mdparamIdDocAll}{38} % Identificativo allegato
\newcommand{\mdparamDesAll}{39} % Descrizione allegato
% Classificazione
\newcommand{\mdparamIndCla}{40} % Indice di classificazione
\newcommand{\mdparamDesCla}{41} % Descrizione (dell'indice di classificazione)
\newcommand{\mdparamPiaCla}{42} % Piano di classificazione
% Riservato
\newcommand{\mdparamRis}{43} % Riservato
% IdentificativoDelFormato
\newcommand{\mdparamFor}{44} % Formato
\newcommand{\mdparamNomPro}{45} % Nome prodotto (software)
\newcommand{\mdparamVerPro}{46} % Versione prodotto (software)
\newcommand{\mdparamProd}{47} % Produttore (software)
% Verifica
\newcommand{\mdparamFirDig}{48} % Firmato digitalmente
\newcommand{\mdparamSigEle}{49} % Sigillato elettronicamente
\newcommand{\mdparamMarTem}{50} % Marcatura temporale
\newcommand{\mdparamConCop}{51} % Conformità copie su supporto informatico
% Agg
% IdAgg
\newcommand{\mdparamTipAgg}{52} % Tipo aggregazione
\newcommand{\mdparamIdAgg}{53} % Identificativo aggregazione
% TipologiaFascicolo
\newcommand{\mdparamTipFas}{54} % Tipologia fascicolo
% NomeDelDocumento
\newcommand{\mdparamNomDoc}{55} % Nome del documento
% VersioneDelDocumento
\newcommand{\mdparamVerDoc}{56} % Versione del documento
% TracciaturaModificheDocumento
\newcommand{\mdparamTipMod}{57} % Tipo modifica
\newcommand{\mdparamDatMod}{58} % Data modifica
\newcommand{\mdparamOraMod}{59} % Ora modifica
% TempoDiConservazione
\newcommand{\mdparamTemCon}{60} % Tempo di conservazione
% Note
\newcommand{\mdparamNot}{61} % Note

\newcommand{\mdsearch}{62}
\newcommand{\searchres}{63}
\newcommand{\nosearchres}{64}
\newcommand{\docpreview}{65}
\newcommand{\nodocpreview}{66}
\newcommand{\exportdoc}{67}
\newcommand{\exportdocfail}{68}
\newcommand{\semparam}{69}
\newcommand{\semsearch}{70}
\newcommand{\printdoc}{71}
\newcommand{\printdocfail}{72}
\newcommand{\verdigsign}{73}
\newcommand{\viewverdigsign}{74}

\begin{document}
    \include{docTitlepage}

    \section*{Registro delle modifiche} {
        \begin{table}[h!]
            \rowcolors{2}{lightgray}{white}
                \begin{tabularx}{\textwidth}{|l|l|l|l|L|L|}
                \hline
                \textbf{Vers.} & \textbf{Data} & \textbf{Autore} & \textbf{Ruolo} & \textbf{Verifica} & \textbf{Oggetto} \\
                \hline
                \textbf{0.15} & 07/02/26 & Francesco Balestro & Analista & Filippo Compagno & Correzioni varie \\
                \hline
                \textbf{0.14} & 07/01/26 & Francesco Balestro & Analista & Mattia Oliva Medin & Inserimento Requisiti di Qualità e di Vincolo e correzioni varie \\
                \hline
                \textbf{0.13} & 03/01/26 & Andrea Masiero & Analista & Filippo Compagno & Inserimento Requisiti Funzionali \\
                \hline
                \textbf{0.12} & 21/12/25 &  Sara Gioia Fichera & Analista & Nenad Radulovic & Inserimento Requisiti Funzionali \\
                \hline
                \textbf{0.11} & 19/12/25 & Bianca Zaghetto & Analista & Nenad Radulovic & Inserimento Casi d'Uso \\
                \hline
                \textbf{0.10} & 17/12/26 & Nenad Radulovic & Analista & Andrea Masiero & Correzioni Casi d'Uso \\
                \hline
                \textbf{0.9} & 16/12/25 & Francesco Balestro & Analista & Andrea Masiero & Inserimento Casi d'Uso \\
                \hline
                \textbf{0.8} & 10/12/25 & Filippo Compagno & Analista & Sara Gioia Fichera & Inserimento Casi d'Uso \\
                \hline
                \textbf{0.7} & 08/12/25 & Mattia Oliva Medin & Analista & Bianca Zaghetto & Inserimento Casi d'Uso \\
                \hline
                \textbf{0.6} & 05/12/25 & Filippo Compagno & Analista & Bianca Zaghetto & Inserimento Casi d'Uso \\
                \hline
                \textbf{0.5} & 04/12/25 & Andrea Masiero & Analista & Sara Gioia Fichera & Inserimento Casi d'Uso \\
                \hline
                \textbf{0.4} & 01/12/25 & Filippo Compagno & Analista & Bianca Zaghetto & Inserimento Casi d'Uso \\
                \hline
                \textbf{0.3} & 24/11/25 & Sara Gioia Fichera & Analista & Francesco Balestro &  Definizione Attori \\
                \hline
                \textbf{0.2} & 20/11/25 & Filippo Compagno & Analista & Francesco Balestro & Caratteristiche degli utenti \\
                \hline
                \textbf{0.1} & 07/11/25 & Sara Gioia Fichera & Analista & Andrea Masiero & Struttura iniziale del documento \\
                \hline
            \end{tabularx}
        \end{table}
    }
    \newpage
    
    \tableofcontents
    \newpage

    \listoffigures
    \newpage
    
    \listoftables
    \newpage

    \section{Introduzione} {
        \subsection{Scopo del documento} {        
            Questo documento ha lo scopo di descrivere in maniera esaustivo e precisa i requisiti funzionali e non funzionali del sistema software da sviluppare per il progetto DIPReader. Per questo motivo integra un'attenta descizione dei casi d'uso, che sono fonte di requisiti.
            Tale analisi è il risultato di un attento studio del capitolato, integrato con quanto concordato durante gli incontri svolti con Sanmarco Informatica, la proponente.
        }
        \subsection{Glossario} {
            Si rimanda al documento Glossario per chiarire i termini nel dominio del progetto che possono risultare ambigui o di difficile comprensione. I vocaboli presenti nel glossario sono stati segnalati seguendo la notazione:
            \begin{center}
                parola\textsubscript{G}
            \end{center}
        }
        \subsection{Riferimenti} {
            \subsubsection{Riferimenti normativi} {
                \begin{itemize}
                    \item Norme di Progetto
                    \item \hypertarget{capitolato}{Capitolato}
                    \begin{itemize}
                        \item {\url{https://www.math.unipd.it/~tullio/IS-1/2025/Progetto/C3.pdf}}
                        \item[] (Ultimo accesso: 2025/11/07)
                    \end{itemize}
                    \item Regolamento del Progetto Didattico
                    \begin{itemize}
                        \item \url{https://www.math.unipd.it/~tullio/IS-1/2025/Dispense/PD1.pdf}
                        \item[] (Ultimo accesso: 2025/11/07)
                    \end{itemize}
                    \item 830-1998 - IEEE Recommended Practice for Software Requirements Specifications
                    \begin{itemize}
                        \item \url{https://ieeexplore.ieee.org/document/720574/references#references}
                        \item[] (Ultimo accesso: )
                    \end{itemize}
                \end{itemize}
            }
            \subsubsection{Riferimenti informativi} {
                \begin{itemize}
                    %\item Slide del corso: T4 - Gestione di progetto
                    %\begin{itemize}
                    %    \item \url{https://www.math.unipd.it/~tullio/IS-1/2025/Dispense/T04.pdf}
                    %    \item[] (Ultimo accesso: 2025/11/07)
                    %\end{itemize}
                    \item Slide del corso: T5 - Analisi dei requisiti
                    \begin{itemize}
                        \item \url{https://www.math.unipd.it/~tullio/IS-1/2025/Dispense/T05.pdf}
                        \item[] (Ultimo accesso: 2025/11/07)
                    \end{itemize}
                    \item Slide del corso: P2 - Diagrammi dei Casi d'Uso
                    \begin{itemize}
                        \item \url{https://www.math.unipd.it/~rcardin/swea/2022/Diagrammi\%20Use\%20Case.pdf}
                        \item[] (Ultimo accesso: 2025/11/07)
                    \end{itemize}
                    \item Repository Github del Prof. Cardin
                    \begin{itemize}
                        \item \url{https://github.com/rcardin/swe-imdb}
                        \item[] (Ultimo accesso: )
                    \end{itemize}
                    \item Glossario
                \end{itemize}
            }
        }
    }

    \newpage
    \section{Descrizione del prodotto} {
        \subsection{Presentazione del prodotto} {
            DIPReader ha lo scopo di supportare la gestione e la conservazione digitale dei documenti, consentendo una transizione strutturata da archivi cartacei a soluzioni cloud avanzate. Il sistema centralizza i documenti e ne assicura autenticità, integrità, leggibilità, reperibilità e sicurezza, offrendo funzionalità di acquisizione, estrazione dei dati e consultazione. Inoltre, permette di richiedere la disponibilità offline dei documenti, garantendo sempre la conformità ai requisiti normativi. 
        }
        \subsection{Funzioni del prodotto} {
            Le funzioni del prodotto sono: 
            \begin{itemize}
                \item Rappresentazione semplificata del contenuto 
                \item Ricerca dei documenti tramite metadati
                \item Visualizzazione in anteprima di alcuni formati
                \item Salvare documenti nel computer dell'Utente
                \item Funzionamento offline
            \end{itemize}
        }
        \subsection{Caratteristiche dell'utente} {
            L’applicazione è destinata a un pubblico ampio e variegato, comprendente ricercatori, professionisti, cittadini interessati o chiunque abbia bisogno di consultare documenti archiviati. Gli utenti possono avere competenze tecniche e digitali molto diverse: dai non esperti, che desiderano trovare rapidamente informazioni, ai professionisti abituati a interpretare metadati complessi. \\
            \\
            Tutti gli utenti condividono l’esigenza di effettuare ricerche efficaci e mirate, utilizzando filtri sui metadati, consultare il contenuto dei documenti e scaricare informazioni in modo affidabile. L’interfaccia dovrà quindi garantire usabilità, chiarezza dei metadati principali e supportare la consultazione e l’estrazione dei documenti in maniera intuitiva.
        }

    }
    \newpage
    \section{Casi d'uso} {
        \subsection{Scopo} {
            Lo scopo di questa sezione è spiegare singolarmente ciascun Caso d'uso individuato.  
        }
        \subsection{Attori} {
            \begin{itemize}
                \item Utente: in questo contesto è l'Attore principale, cioè colui che intergisce attivamente con il sistema e svolge le azioni indicate dal Caso d'Uso.
                \item Motore semantico: si occupa di elaborare una ricerca semantica fatta dall'utente per produrre un risultato. In questo contesto è un Attore secondario, in quanto interagisce passivamente con il sistema.
                \item Gestore di stampa SO: si occupa di gestire la richiesta di stampa di un documento da parte dell'utente. In questo contesto è un Attore secondario, in quanto interagisce passivamente con il sistema.
            \end{itemize}
        }
        \subsection{Casi d'uso} {
            % Caricamento DIP
            \begin{usecase}
                \code{\loaddip}
                \titleUC{Caricamento DIP}
                \actor{Utente}
                \pre{
                    \item Nessun DIP è caricato nel sistema
                }
                \post{
                    \item Il DIP è caricato nel sistema
                }
                \main{
                    \item L'utente seleziona la funzione di caricamento
                    \item L'utente seleziona il pacchetto (\textbf{UC\loaddip.1})
                    \item L'utente visualizza la conferma dell'avvenuto caricamento
                }
                \alt{
                    \item[2a.] L'utente cancella l'operazione
                    \item[3a.] L'utente inserisce un pacchetto con struttura non conforme (\textbf{UC\errloaddip})
                }
                \inc{\ucref{uc:\loaddip.1}}
                \ext{\ucref{uc:\errloaddip}}
                \gen{\ucref{uc:\remloaddip}}
            \end{usecase}
            \begin{figure}[H]
                    \centering
                    \includegraphics[width=0.6\textwidth]{Use_Case_Diagrams/UC01.png}
                    \caption{UC01 - Caricamento DIP}
            \end{figure}
            
            \begin{usecase}
                \code{\loaddip.1}
                \titleUC{Selezione pacchetto}
                \actor{Utente}
                \pre{
                    \item L'utente ha selezionato la funzione di caricamento
                }
                \post{
                    \item L'utente ha selezionato la funzione di caricamento
                }
                \main{
                    \item Il sistema richiede il pacchetto all'utente
                    \item L'utente seleziona il pacchetto dal proprio dispositivo
                }
            \end{usecase}
            
            % Caricamento DIP con DIP già caricato
            \begin{usecase}
                \code{\remloaddip}
                \titleUC{Caricamento DIP con DIP già presente}
                \actor{Utente}
                \pre{
                    \item Un DIP è caricato nel sistema
                }
                \post{
                    \item Il DIP già caricato è rimosso dal sistema
                    \item Il nuovo DIP è caricato nel sistema
                }
                \main{
                    \item L'utente seleziona la funzione di caricamento
                    \item L'utente vede un messaggio informativo sulla necessità di rimozione del DIP presente 
                    \item L'utente rimuove il DIP presente nel sistema (\textbf{UC\removedip})
                    \item L'utente seleziona il pacchetto (\textbf{UC\loaddip.1})
                    \item L'utente visualizza la conferma dell'avvenuto caricamento
                }
                \alt{
                    \item[4a.] L'utente cancella l'operazione
                    \item[5a.] L'utente inserisce un pacchetto con struttura non conforme (\textbf{UC\errloaddip})
                }
                \inc{\ucref{uc:\removedip}}
            \end{usecase}

            % Errore struttura DIP
            \begin{usecase}
                \code{\errloaddip}
                \titleUC{Errore struttura DIP}
                \actor{Utente}
                \pre{
                    \item L'utente ha caricato un pacchetto che non rappresenta un DIP
                }
                \post{
                    \item Il DIP non viene caricato nel sistema
                    \item L'utente vede un messaggio di errore
                }
                \main{
                    \item L'utente vede un messaggio di errore relativo alla struttura non conforme del DIP
                }
            \end{usecase}
         
            % Rimozione DIP
            \begin{usecase}
                \code{\removedip}
                \titleUC{Rimozione DIP}
                \actor{Utente}
                \pre{
                    \item L'utente ha caricato un DIP nel sistema
                }
                \post{
                    \item Il DIP caricato è rimosso dal sistema
                }
                \main{
                    \item L'utente seleziona la funzionalità di chiusura
                    \item Il sistema richiede la conferma
                    \item L'utente conferma la chiusura del DIP
                }
                \alt{
                    \item[3a.] L'utente cancella l'operazione
                }
            \end{usecase}
            
            % Visualizzazione DIP
            \begin{usecase}
                \code{\viewdip}
                \titleUC{Visualizzazione contenuto DIP}
                \actor{Utente}
                \pre{
                    \item L'utente ha caricato un DIP nel sistema
                }
                \post{
                    \item L'utente vede le informazioni relative al DIP
                    \item L'utente vede i documenti contenuti nel DIP
                }
                \main{
                    \item L'utente visualizza le informazioni relative al DIP (\textbf{UC\viewdip.1})
                    \item L'utente visualizza la lista dei documenti contenuti nel DIP (\textbf{UC\viewdip.2})
                }
                \inc{\ucref{uc:\viewdip.1}, \ucref{uc:\viewdip.2}}
            \end{usecase}
            \begin{figure}[H]
                    \centering
                    \includegraphics[width=0.6\textwidth]{Use_Case_Diagrams/UC05.png}
                    \caption{UC05 - Visualizzazione contenuto DIP}
            \end{figure}
            \begin{usecase}
                \code{\viewdip.1}
                \titleUC{Visualizzazione informazioni DIP}
                \actor{Utente}
                \pre{
                    \item L'utente sta visualizzando il contenuto del DIP
                }
                \post{
                    \item L'utente vede le informazioni relative al DIP
                }
                \main{
                    \item L'utente visualizza
                    \begin{itemize}
                        \item La data di creazione del DIP
                        \item Le informazioni del programma usato per la creazione del DIP, quali
                        \begin{itemize}
                            \item Nome
                            \item Versione
                            \item Produttore
                            \item Il numero di documenti contenuti
                        \end{itemize}
                    \end{itemize}
                }
            \end{usecase}
            \begin{usecase}
                \code{\viewdip.2}
                \titleUC{Visualizzazione lista documenti contenuti}
                \actor{Utente}
                \pre{
                    \item L'utente sta visualizzando il contenuto del DIP
                }
                \post{
                    \item L'utente vede la lista di documenti contenuti nel DIP
                }
                \main{
                    \item L'utente visualizza i documenti contenuti nel DIP
                }
                \inc{\ucref{uc:\viewdip.2.1}}
            \end{usecase}
            \begin{usecase}
                \code{\viewdip.2.1}
                \titleUC{Visualizzazione singolo documento}
                \actor{Utente}
                \pre{
                    \item L'utente sta visualizzando la lista dei documenti contenuti nel DIP
                }
                \post{
                    \item L'utente vede le informazioni principali di un documento contenuto nel DIP
                }
                \main{
                    \item L'utente visualizza
                    \begin{itemize}
                        \item Il nome del file
                        \item Il suo formato
                        \item Identificativo documentale
                        \item Data in cui il documento è stato registrato
                    \end{itemize}
                }
                \inc{\ucref{uc:\viewdetdoc}}
                \gen{\ucref{uc:\viewdip.2.2}}
            \end{usecase}
            \begin{usecase}
                \code{\viewdip.2.2}
                \titleUC{Visualizzazione singolo documento con allegati}
                \actor{Utente}
                \pre{
                    \item L'utente sta visualizzando la lista dei documenti contenuti nel DIP
                }
                \post{
                    \item L'utente vede le informazioni principali di un documento contenuto nel DIP
                    \item L'utente vede la lista degli allegati del documento
                }
                \main{
                    \item L'utente visualizza
                    \begin{itemize}
                        \item Il nome del file
                        \item Il suo formato
                        \item Identificativo documentale
                        \item Data in cui il documento è stato registrato
                        \item La lista dei suoi allegati (\textbf{UC\viewdip.2.2.1})
                    \end{itemize}
                }
                \inc{\ucref{uc:\viewdip.2.2.1}}
            \end{usecase}
            \begin{usecase}
                \code{\viewdip.2.2.1}
                \titleUC{Visualizzazione lista allegati documento}
                \actor{Utente}
                \pre{
                    \item L'utente sta visualizzando un documento contenuto nel DIP
                }
                \post{
                    \item L'utente vede la lista degli allegati di un documento contenuto nel DIP
                }
                \main{
                    \item L'utente visualizza gli allegati di un documento contenuto nel DIP
                }
                \inc{\ucref{uc:\viewdip.2.2.1.1}}
            \end{usecase}
            \begin{usecase}
                \code{\viewdip.2.2.1.1}
                \titleUC{Visualizzazione singolo allegato}
                \actor{Utente}
                \pre{
                    \item L'utente sta visualizzando la lista degli allegati di un documento contenuto nel DIP
                }
                \post{
                    \item L'utente vede le informazioni principali di un allegato di un documento contenuto nel DIP
                }
                \main{
                    \item L'utente visualizza
                    \begin{itemize}
                        \item Il nome dell'allegato
                        \item Il suo formato
                        \item Identificativo allegato
                    \end{itemize}
                }
            \end{usecase}

            % Visualizzazione dettagli documento
            \begin{usecase}
                \code{\viewdetdoc}
                \titleUC{Visualizzazione dettagli documento}
                \actor{Utente}
                \pre{
                    \item L'utente sta visualizzando un singolo documento
                }
                \post{
                    \item L'utente vede i dettagli di un documento
                }
                \main{
                    \item L'utente seleziona la funzionalità di visualizzazione dei dettagli di un documento
                    \item L'utente visualizza i dettagli di un documento
                }
            \end{usecase}
            
            % Visualizzazione delle informazioni di conservazione
            \begin{usecase}
                \code{\viewconsinfo}
                \titleUC{Visualizzazione informazioni di conservazione}
                \actor{Utente}
                \pre{
                    \item L'utente ha caricato un DIP
                    \item L'utente sta visualizzando una lista di documenti contenuti nel DIP
                }
                \post{
                    \item L’utente visualizza i dettagli relativi alla validità legale e allo stato del processo di conservazione dei documenti
                }
                \main{
                    \item L'utente visualizza degli elementi grafici che gli indicano la validità e lo stato del processo di conservazione dei documenti
                }
            \end{usecase}
            \begin{figure}[H]
                    \centering
                    \includegraphics[width=0.5\textwidth]{Use_Case_Diagrams/UC07.png}
                    \caption{UC07 - Visualizzazione informazioni di conservazione}
            \end{figure}

            % Criteri di ricerca (metadati)
            \begin{usecase}
                \code{\mdparamImpCrit}
                \titleUC{Inserimento criterio di ricerca "Impronta crittografica"}
                \actor{Utente}
                \pre{
                    \item L'utente ha caricato un DIP
                    \item L'utente sta visualizzando una lista di documenti contenuti nel DIP
                }
                \post{
                    \item L'utente ha configurato un criterio di ricerca
                }
                \main{
                    \item L'utente seleziona la funzionalità di aggiunta criterio di ricerca
                    \item L'utente seleziona il criterio "Impronta crittografica"
                    \item L'utente inserisce l'hash del documento (\textbf{UC\mdparamImpCrit.1})
                    \item L'utente seleziona l'algoritmo applicato (\textbf{UC\mdparamImpCrit.2})
                }
                \inc{\ucref{uc:\mdparamImpCrit.1}, \ucref{uc:\mdparamImpCrit.2}}
                \gen{\ucref{uc:\mdparamImpCritAll}}
            \end{usecase}
            \begin{figure}[H]
                    \centering
                    \includegraphics[width=0.6\textwidth]{Use_Case_Diagrams/UC08.png}
                    \caption{UC08 - Inserimento criterio di ricerca ”Impronta crittografica”}
            \end{figure}
            \begin{usecase}
                \code{\mdparamImpCrit.1}
                \titleUC{Inserimento hash del documento}
                \actor{Utente}
                \pre{
                    \item L'utente sta configurando il criterio di ricerca "Impronta crittografica"
                }
                \post{
                    \item L'utente ha inserito l'hash del documento
                }
                \main{
                    \item L'utente inserisce l'hash del documento
                }
            \end{usecase}
            \begin{usecase}
                \code{\mdparamImpCrit.2}
                \titleUC{Selezione algoritmo di hashing}
                \actor{Utente}
                \pre{
                    \item L'utente sta configurando il criterio di ricerca "Impronta crittografica"
                }
                \post{
                    \item L'utente ha selezionato l'algoritmo di hashing
                }
                \main{
                    \item L'utente seleziona l'algoritmo di hashing applicato tra quelli disponibili
                }
            \end{usecase}

            \begin{usecase}
                \code{\mdparamIdDoc}
                \titleUC{Inserimento criterio di ricerca "Identificativo documento"}
                \actor{Utente}
                \pre{
                    \item L'utente ha caricato un DIP
                    \item L'utente sta visualizzando una lista di documenti contenuti nel DIP
                }
                \post{
                    \item L'utente ha configurato un criterio di ricerca
                }
                \main{
                    \item L'utente seleziona la funzionalità di aggiunta criterio di ricerca
                    \item L'utente seleziona il criterio "Identificativo documento"
                    \item L'utente inserisce il valore
                }
                \gen{\ucref{uc:\mdparamIdDocAll}}
            \end{usecase}
            \begin{figure}[H]
                    \centering
                    \includegraphics[width=0.5\textwidth]{Use_Case_Diagrams/UC09.png}
                    \caption{UC09 - Inserimento criterio di ricerca ”Identificativo documento”}
            \end{figure}

            \begin{usecase}
                \code{\mdparamModFor}
                \titleUC{Inserimento criterio di ricerca "Modalità di formazione"}
                \actor{Utente}
                \pre{
                    \item L'utente ha caricato un DIP
                    \item L'utente sta visualizzando una lista di documenti contenuti nel DIP
                }
                \post{
                    \item L'utente ha configurato un criterio di ricerca
                }
                \main{
                    \item L'utente seleziona la funzionalità di aggiunta criterio di ricerca
                    \item L'utente seleziona il criterio "Modalità di formazione"
                    \item L'utente seleziona i valori tra le opzioni
                    \begin{itemize}
                        \item "Creazione tramite l'utilizzo di strumenti software che assicurino la produzione di documenti nei formati previsti nell'Allegato 2 delle Linee Guida"
                        \item "Acquisizione di un documento informatico per via telematica o su supporto informatico, acquisizione della copia per immagine su supporto informatico di un documento analogico, acquisizione della copia informatica di un documento analogico"
                        \item "Memorizzazione su supporto informatico in formato digitale delle informazioni risultanti da transazioni o processi informatici o dalla presentazione telematica di dati attraverso moduli o formulari resi disponibili all'utente"
                        \item "Generazione o raggruppamento anche in via automatica di un insieme di dati o registrazioni, provenienti da una o più banche dati, anche appartenenti a più soggetti interoperanti, secondo una struttura logica predeterminata e memorizzata in forma statica"
                    \end{itemize}
                }
            \end{usecase}
            \begin{figure}[H]
                    \centering
                    \includegraphics[width=0.5\textwidth]{Use_Case_Diagrams/UC10.png}
                    \caption{UC10 - Inserimento criterio di ricerca ”Modalità di formazione”}
            \end{figure}

            \begin{usecase}
                \code{\mdparamTipDoc}
                \titleUC{Inserimento criterio di ricerca "Tipologia documentale"}
                \actor{Utente}
                \pre{
                    \item L'utente ha caricato un DIP
                    \item L'utente sta visualizzando una lista di documenti contenuti nel DIP
                }
                \post{
                    \item L'utente ha configurato un criterio di ricerca
                }
                \main{
                    \item L'utente seleziona la funzionalità di aggiunta criterio di ricerca
                    \item L'utente seleziona il criterio "Tipologia documentale"
                    \item L'utente inserisce il valore
                }
            \end{usecase}
            \begin{figure}[H]
                    \centering
                    \includegraphics[width=0.5\textwidth]{Use_Case_Diagrams/UC11.png}
                    \caption{UC11 - Inserimento criterio di ricerca ”Tipologia documentale”}
            \end{figure}

            \begin{usecase}
                \code{\mdparamTipFlu}
                \titleUC{Inserimento criterio di ricerca "Tipologia di flusso"}
                \actor{Utente}
                \pre{
                    \item L'utente ha caricato un DIP
                    \item L'utente sta visualizzando una lista di documenti contenuti nel DIP
                }
                \post{
                    \item L'utente ha configurato un criterio di ricerca
                }
                \main{
                    \item L'utente seleziona la funzionalità di aggiunta criterio di ricerca
                    \item L'utente seleziona il criterio "Tipologia di flusso"
                    \item L'utente seleziona i valori tra le opzioni
                    \begin{itemize}
                        \item "U (In uscita)"
                        \item "E (In entrata)"
                        \item "I (Interno)"
                    \end{itemize}
                }
            \end{usecase}
            \begin{figure}[H]
                    \centering
                    \includegraphics[width=0.5\textwidth]{Use_Case_Diagrams/UC12.png}
                    \caption{UC12 - Inserimento criterio di ricerca ”Tipologia di flusso”}
            \end{figure}

            \begin{usecase}
                \code{\mdparamTipReg}
                \titleUC{Inserimento criterio di ricerca "Tipo registro"}
                \actor{Utente}
                \pre{
                    \item L'utente ha caricato un DIP
                    \item L'utente sta visualizzando una lista di documenti contenuti nel DIP
                }
                \post{
                    \item L'utente ha configurato un criterio di ricerca
                }
                \main{
                    \item L'utente seleziona la funzionalità di aggiunta criterio di ricerca
                    \item L'utente seleziona il criterio "Tipo registro"
                    \item L'utente seleziona il valore tra le opzioni
                    \begin{itemize}
                        \item "Nessuno"
                        \item "Protocollo ordinario/Protocollo emergenza"
                        \item "Repertorio/Registro"
                    \end{itemize}
                }
            \end{usecase}
            \begin{figure}[H]
                    \centering
                    \includegraphics[width=0.5\textwidth]{Use_Case_Diagrams/UC13.png}
                    \caption{UC13 - Inserimento criterio di ricerca ”Tipo registro”}
            \end{figure}

            \begin{usecase}
                \code{\mdparamDatReg}
                \titleUC{Inserimento criterio di ricerca "Data registrazione"}
                \actor{Utente}
                \pre{
                    \item L'utente ha caricato un DIP
                    \item L'utente sta visualizzando una lista di documenti contenuti nel DIP
                }
                \post{
                    \item L'utente ha configurato un criterio di ricerca
                }
                \main{
                    \item L'utente seleziona la funzionalità di aggiunta criterio di ricerca
                    \item L'utente seleziona il criterio "Data registrazione"
                    \item L'utente inserisce il valore
                }
            \end{usecase}
            \begin{figure}[H]
                    \centering
                    \includegraphics[width=0.5\textwidth]{Use_Case_Diagrams/UC14.png}
                    \caption{UC14 - Inserimento criterio di ricerca ”Data registrazione”}
            \end{figure}

            \begin{usecase}
                \code{\mdparamOraReg}
                \titleUC{Inserimento criterio di ricerca "Ora registrazione"}
                \actor{Utente}
                \pre{
                    \item L'utente ha caricato un DIP
                    \item L'utente sta visualizzando una lista di documenti contenuti nel DIP
                }
                \post{
                    \item L'utente ha configurato un criterio di ricerca
                }
                \main{
                    \item L'utente seleziona la funzionalità di aggiunta criterio di ricerca
                    \item L'utente seleziona il criterio "Ora registrazione"
                    \item L'utente inserisce il valore
                }
            \end{usecase}
            \begin{figure}[H]
                    \centering
                    \includegraphics[width=0.5\textwidth]{Use_Case_Diagrams/UC15.png}
                    \caption{UC15 - Inserimento criterio di ricerca ”Ora registrazione”}
            \end{figure}

            \begin{usecase}
                \code{\mdparamNumDoc}
                \titleUC{Inserimento criterio di ricerca "Numero documento"}
                \actor{Utente}
                \pre{
                    \item L'utente ha caricato un DIP
                    \item L'utente sta visualizzando una lista di documenti contenuti nel DIP
                }
                \post{
                    \item L'utente ha configurato un criterio di ricerca
                }
                \main{
                    \item L'utente seleziona la funzionalità di aggiunta criterio di ricerca
                    \item L'utente seleziona il criterio "Numero documento"
                    \item L'utente inserisce il valore
                }
            \end{usecase}
            \begin{figure}[H]
                    \centering
                    \includegraphics[width=0.5\textwidth]{Use_Case_Diagrams/UC16.png}
                    \caption{UC16 - Inserimento criterio di ricerca ”Numero documento”}
            \end{figure}

            \begin{usecase}
                \code{\mdparamCodReg}
                \titleUC{Inserimento criterio di ricerca "Codice registro"}
                \actor{Utente}
                \pre{
                    \item L'utente ha caricato un DIP
                    \item L'utente sta visualizzando una lista di documenti contenuti nel DIP
                }
                \post{
                    \item L'utente ha configurato un criterio di ricerca
                }
                \main{
                    \item L'utente seleziona la funzionalità di aggiunta criterio di ricerca
                    \item L'utente seleziona il criterio "Codice registro"
                    \item L'utente inserisce il valore
                }
            \end{usecase}
            \begin{figure}[H]
                    \centering
                    \includegraphics[width=0.5\textwidth]{Use_Case_Diagrams/UC17.png}
                    \caption{UC17 - Inserimento criterio di ricerca ”Codice registro”}
            \end{figure}

             \begin{usecase}
                \code{\mdparamRuolo}
                \titleUC{Inserimento criterio di ricerca "Ruolo"}
                \actor{Utente}
                \pre{
                    \item L'utente ha caricato un DIP
                    \item L'utente sta visualizzando una lista di documenti contenuti nel DIP
                }
                \post{
                    \item L'utente ha configurato un criterio di ricerca
                }
                \main{
                    \item L'utente seleziona la funzionalità di aggiunta criterio di ricerca
                    \item L'utente seleziona il criterio "Ruolo"
                    \item L'utente seleziona il valore tra le opzioni
                    \begin{itemize}
                        \item "Altro"
                        \item "Assegnatario"
                        \item "Autore"
                        \item "Destinatario"
                        \item "Mittente"
                        \item "Operatore"
                        \item "Produttore"
                        \item "Responsabile del Servizio Protocollo"
                        \item "Responsabile della Gestione Documentale"
                        \item "Soggetto che effettua la registrazione"
                    \end{itemize}
                }
            \end{usecase}
            \begin{figure}[H]
                    \centering
                    \includegraphics[width=0.5\textwidth]{Use_Case_Diagrams/UC18.png}
                    \caption{UC18 - Inserimento criterio di ricerca ”Ruolo”}
            \end{figure}

            \begin{usecase}
                \code{\mdparamTipSog}
                \titleUC{Inserimento criterio di ricerca "Tipo soggetto"}
                \actor{Utente}
                \pre{
                    \item L'utente ha caricato un DIP
                    \item L'utente sta visualizzando una lista di documenti contenuti nel DIP
                }
                \post{
                    \item L'utente ha configurato un criterio di ricerca
                }
                \main{
                    \item L'utente seleziona la funzionalità di aggiunta criterio di ricerca
                    \item L'utente seleziona il criterio "Tipo soggetto"
                    \item L'utente seleziona il valore tra le opzioni disponibili
                    \begin{itemize}
                        \item "AS (Assegnatario)"
                        \item "PAE (Amministrazione pubblica estera)"
                        \item "PAI (Amministrazione pubblica italiana)"
                        \item "PF (Persona Fisica)"
                        \item "PG (Organizzazione)"
                        \item "SW (Documenti prodotti automaticamente)"
                    \end{itemize}
                }
            \end{usecase}
            \begin{figure}[H]
                    \centering
                    \includegraphics[width=0.5\textwidth]{Use_Case_Diagrams/UC19.png}
                    \caption{UC19 - Inserimento criterio di ricerca ”Tipo soggetto”}
            \end{figure}

            \begin{usecase}
                \code{\mdparamCog}
                \titleUC{Inserimento criterio di ricerca "Cognome"}
                \actor{Utente}
                \pre{
                    \item L'utente ha caricato un DIP
                    \item L'utente sta visualizzando una lista di documenti contenuti nel DIP
                    \item L'utente ha configurato il criterio di ricerca "Ruolo" con un valore tra
                    \begin{itemize}
                        \item "Altro"
                        \item "Assegnatario"
                        \item "Autore"
                        \item "Destinatario"
                        \item "Mittente"
                        \item "Operatore"
                        \item "Responsabile del Servizio Protocollo"
                        \item "Responsabile della Gestione Documentale"
                        \item "Soggetto che effettua la registrazione"
                    \end{itemize}
                    \item L'utente ha configurato il criterio di ricerca "Tipo soggetto" con un valore tra
                    \begin{itemize}
                        \item "AS (Assegnatario)"
                        \item "PF (Persona Fisica)"
                    \end{itemize}
                }
                \post{
                    \item L'utente ha configurato un criterio di ricerca
                }
                \main{
                    \item L'utente seleziona la funzionalità di aggiunta criterio di ricerca
                    \item L'utente seleziona il criterio "Cognome"
                    \item L'utente inserisce il valore
                }
            \end{usecase}
            \begin{figure}[H]
                    \centering
                    \includegraphics[width=0.5\textwidth]{Use_Case_Diagrams/UC20.png}
                    \caption{UC20 - Inserimento criterio di ricerca ”Cognome”}
            \end{figure}

            \begin{usecase}
                \code{\mdparamNom}
                \titleUC{Inserimento criterio di ricerca "Nome"}
                \actor{Utente}
                \pre{
                    \item L'utente ha caricato un DIP
                    \item L'utente sta visualizzando una lista di documenti contenuti nel DIP
                    \item L'utente ha configurato il criterio di ricerca "Ruolo" con un valore tra
                    \begin{itemize}
                        \item "Altro"
                        \item "Assegnatario"
                        \item "Autore"
                        \item "Destinatario"
                        \item "Mittente"
                        \item "Operatore"
                        \item "Responsabile del Servizio Protocollo"
                        \item "Responsabile della Gestione Documentale"
                        \item "Soggetto che effettua la registrazione"
                    \end{itemize}
                    \item L'utente ha configurato il criterio di ricerca "Tipo soggetto" con un valore tra
                    \begin{itemize}
                        \item "AS (Assegnatario)"
                        \item "PF (Persona Fisica)"
                    \end{itemize}
                }
                \post{
                    \item L'utente ha configurato un criterio di ricerca
                }
                \main{
                    \item L'utente seleziona la funzionalità di aggiunta criterio di ricerca
                    \item L'utente seleziona il criterio "Nome"
                    \item L'utente inserisce il valore
                }
            \end{usecase}
            \begin{figure}[H]
                    \centering
                    \includegraphics[width=0.5\textwidth]{Use_Case_Diagrams/UC21.png}
                    \caption{UC21 - Inserimento criterio di ricerca ”Nome”}
            \end{figure}

            \begin{usecase}
                \code{\mdparamCF}
                \titleUC{Inserimento criterio di ricerca "Codice Fiscale"}
                \actor{Utente}
                \pre{
                    \item L'utente ha caricato un DIP
                    \item L'utente sta visualizzando una lista di documenti contenuti nel DIP
                    \item L'utente ha configurato il criterio di ricerca "Ruolo" con un valore tra
                    \begin{itemize}
                        \item "Altro"
                        \item "Assegnatario"
                        \item "Autore"
                        \item "Destinatario"
                        \item "Mittente"
                        \item "Operatore"
                        \item "Responsabile del Servizio Protocollo"
                        \item "Responsabile della Gestione Documentale"
                        \item "Soggetto che effettua la registrazione"
                    \end{itemize}
                    \item L'utente ha configurato il criterio di ricerca "Tipo soggetto" con un valore tra
                    \begin{itemize}
                        \item "AS (Assegnatario)"
                        \item "PF (Persona Fisica)"
                        \item "PG (Organizzazione)"
                    \end{itemize}
                }
                \post{
                    \item L'utente ha configurato un criterio di ricerca
                }
                \main{
                    \item L'utente seleziona la funzionalità di aggiunta criterio di ricerca
                    \item L'utente seleziona il criterio "Codice Fiscale"
                    \item L'utente inserisce il valore
                }
            \end{usecase}
            \begin{figure}[H]
                    \centering
                    \includegraphics[width=0.5\textwidth]{Use_Case_Diagrams/UC22.png}
                    \caption{UC22 - Inserimento criterio di ricerca ”Codice Fiscale”}
            \end{figure}

            \begin{usecase}
                \code{\mdparamIndDig}
                \titleUC{Inserimento criterio di ricerca "Indirizzo digitale"}
                \actor{Utente}
                \pre{
                    \item L'utente ha caricato un DIP
                    \item L'utente sta visualizzando una lista di documenti contenuti nel DIP
                    \item L'utente ha configurato il criterio di ricerca "Ruolo" con un valore tra
                    \begin{itemize}
                        \item "Altro"
                        \item "Assegnatario"
                        \item "Autore"
                        \item "Destinatario"
                        \item "Mittente"
                        \item "Operatore"
                        \item "Responsabile del Servizio Protocollo"
                        \item "Responsabile della Gestione Documentale"
                        \item "Soggetto che effettua la registrazione"
                    \end{itemize}
                    \item L'utente ha configurato il criterio di ricerca "Tipo soggetto" con un valore tra
                    \begin{itemize}
                        \item "AS (Assegnatario)"
                        \item "PAE (Amministrazione pubblica estera)"
                        \item "PAI (Amministrazione pubblica italiana)"
                        \item "PF (Persona Fisica)"
                        \item "PG (Organizzazione)"
                    \end{itemize}
                }
                \post{
                    \item L'utente ha configurato un criterio di ricerca
                }
                \main{
                    \item L'utente seleziona la funzionalità di aggiunta criterio di ricerca
                    \item L'utente seleziona il criterio "Indirizzo digitale"
                    \item L'utente inserisce il valore
                }
            \end{usecase}
            \begin{figure}[H]
                    \centering
                    \includegraphics[width=0.5\textwidth]{Use_Case_Diagrams/UC23.png}
                    \caption{UC23 - Inserimento criterio di ricerca ”Indirizzo digitale”}
            \end{figure}

            \begin{usecase}
                \code{\mdparamDenOrg}
                \titleUC{Inserimento criterio di ricerca "Denominazione organizzazione"}
                \actor{Utente}
                \pre{
                    \item L'utente ha caricato un DIP
                    \item L'utente sta visualizzando una lista di documenti contenuti nel DIP
                    \item L'utente ha configurato il criterio di ricerca "Ruolo" con un valore tra
                    \begin{itemize}
                        \item "Altro"
                        \item "Assegnatario"
                        \item "Autore"
                        \item "Destinatario"
                        \item "Mittente"
                        \item "Soggetto che effettua la registrazione"
                    \end{itemize}
                    \item L'utente ha configurato il criterio di ricerca "Tipo soggetto" con un valore tra
                    \begin{itemize}
                        \item "AS (Assegnatario)"
                        \item "PG (Organizzazione)"
                    \end{itemize}
                }
                \post{
                    \item L'utente ha configurato un criterio di ricerca
                }
                \main{
                    \item L'utente seleziona la funzionalità di aggiunta criterio di ricerca
                    \item L'utente seleziona il criterio "Denominazione organizzazione"
                    \item L'utente inserisce il valore
                }
            \end{usecase}
            \begin{figure}[H]
                    \centering
                    \includegraphics[width=0.5\textwidth]{Use_Case_Diagrams/UC24.png}
                    \caption{UC24 - Inserimento criterio di ricerca ”Denominazione organizzazione”}
            \end{figure}

            \begin{usecase}
                \code{\mdparamCFPI}
                \titleUC{Inserimento criterio di ricerca "Partita IVA"}
                \actor{Utente}
                \pre{
                    \item L'utente ha caricato un DIP
                    \item L'utente sta visualizzando una lista di documenti contenuti nel DIP
                    \item L'utente ha configurato il criterio di ricerca "Ruolo" con un valore tra
                    \begin{itemize}
                        \item "Altro"
                        \item "Autore"
                        \item "Destinatario"
                        \item "Mittente"
                        \item "Soggetto che effettua la registrazione"
                    \end{itemize}
                    \item L'utente ha configurato il criterio di ricerca "Tipo soggetto" con un valore tra
                    \begin{itemize}
                        \item "PG (Organizzazione)"
                    \end{itemize}
                }
                \post{
                    \item L'utente ha configurato un criterio di ricerca
                }
                \main{
                    \item L'utente seleziona la funzionalità di aggiunta criterio di ricerca
                    \item L'utente seleziona il criterio "Partita IVA"
                    \item L'utente inserisce il valore
                }
            \end{usecase}
            \begin{figure}[H]
                    \centering
                    \includegraphics[width=0.5\textwidth]{Use_Case_Diagrams/UC25.png}
                    \caption{UC25 - Inserimento criterio di ricerca ”Partita IVA”}
            \end{figure}

            \begin{usecase}
                \code{\mdparamDenUff}
                \titleUC{Inserimento criterio di ricerca "Denominazione ufficio"}
                \actor{Utente}
                \pre{
                    \item L'utente ha caricato un DIP
                    \item L'utente sta visualizzando una lista di documenti contenuti nel DIP
                    \item L'utente ha configurato il criterio di ricerca "Ruolo" con un valore tra
                    \begin{itemize}
                        \item "Altro"
                        \item "Assegnatario"
                        \item "Autore"
                        \item "Destinatario"
                        \item "Mittente"
                        \item "Soggetto che effettua la registrazione"
                    \end{itemize}
                    \item L'utente ha configurato il criterio di ricerca "Tipo soggetto" con un valore tra
                    \begin{itemize}
                        \item "AS (Assegnatario)"
                        \item "PAE (Amministrazione pubblica estera)"
                        \item "PG (Organizzazione)"
                    \end{itemize}
                }
                \post{
                    \item L'utente ha configurato un criterio di ricerca
                }
                \main{
                    \item L'utente seleziona la funzionalità di aggiunta criterio di ricerca
                    \item L'utente seleziona il criterio "Denominazione ufficio"
                    \item L'utente inserisce il valore
                }
            \end{usecase}
            \begin{figure}[H]
                    \centering
                    \includegraphics[width=0.5\textwidth]{Use_Case_Diagrams/UC26.png}
                    \caption{UC26 - Inserimento criterio di ricerca ”Denominazione ufficio”}
            \end{figure}

            \begin{usecase}
                \code{\mdparamDenAmm}
                \titleUC{Inserimento criterio di ricerca "Denominazione amministrazione"}
                \actor{Utente}
                \pre{
                    \item L'utente ha caricato un DIP
                    \item L'utente sta visualizzando una lista di documenti contenuti nel DIP
                    \item L'utente ha configurato il criterio di ricerca "Ruolo" con un valore tra
                    \begin{itemize}
                        \item "Altro"
                        \item "Autore"
                        \item "Destinatario"
                        \item "Mittente"
                    \end{itemize}
                    \item L'utente ha configurato il criterio di ricerca "Tipo soggetto" con un valore tra
                    \begin{itemize}
                        \item "PAE (Amministrazione pubblica estera)"
                        \item "PAI (Amministrazione pubblica italiana)"
                    \end{itemize}
                }
                \post{
                    \item L'utente ha configurato un criterio di ricerca
                }
                \main{
                    \item L'utente seleziona la funzionalità di aggiunta criterio di ricerca
                    \item L'utente seleziona il criterio "Denominazione amministrazione"
                    \item L'utente inserisce il valore
                }
            \end{usecase}
            \begin{figure}[H]
                    \centering
                    \includegraphics[width=0.5\textwidth]{Use_Case_Diagrams/UC27.png}
                    \caption{UC27 - Inserimento criterio di ricerca ”Denominazione amministrazione”}
            \end{figure}

            \begin{usecase}
                \code{\mdparamIPA}
                \titleUC{Inserimento criterio di ricerca "Codice IPA"}
                \actor{Utente}
                \pre{
                    \item L'utente ha caricato un DIP
                    \item L'utente sta visualizzando una lista di documenti contenuti nel DIP
                    \item L'utente ha configurato il criterio di ricerca "Ruolo" con un valore tra
                    \begin{itemize}
                        \item "Altro"
                        \item "Autore"
                        \item "Destinatario"
                        \item "Mittente"
                    \end{itemize}
                    \item L'utente ha configurato il criterio di ricerca "Tipo soggetto" con il valore "PAI (Amministrazione pubblica italiana)"
                }
                \post{
                    \item L'utente ha configurato un criterio di ricerca
                }
                \main{
                    \item L'utente seleziona la funzionalità di aggiunta criterio di ricerca
                    \item L'utente seleziona il criterio "Codice IPA"
                    \item L'utente inserisce il valore
                }
            \end{usecase}
            \begin{figure}[H]
                    \centering
                    \includegraphics[width=0.5\textwidth]{Use_Case_Diagrams/UC28.png}
                    \caption{UC28 - Inserimento criterio di ricerca ”Codice IPA”}
            \end{figure}

            \begin{usecase}
                \code{\mdparamDenAmmAOO}
                \titleUC{Inserimento criterio di ricerca "Denominazione amministrazione AOO"}
                \actor{Utente}
                \pre{
                    \item L'utente ha caricato un DIP
                    \item L'utente sta visualizzando una lista di documenti contenuti nel DIP
                    \item L'utente ha configurato il criterio di ricerca "Ruolo" con un valore tra
                    \begin{itemize}
                        \item "Altro"
                        \item "Autore"
                        \item "Destinatario"
                        \item "Mittente"
                    \end{itemize}
                    \item L'utente ha configurato il criterio di ricerca "Tipo soggetto" con il valore "PAI (Amministrazione pubblica italiana)"
                }
                \post{
                    \item L'utente ha configurato un criterio di ricerca
                }
                \main{
                    \item L'utente seleziona la funzionalità di aggiunta criterio di ricerca
                    \item L'utente seleziona il criterio "Denominazione amministrazione AOO"
                    \item L'utente inserisce il valore
                }
            \end{usecase}
            \begin{figure}[H]
                    \centering
                    \includegraphics[width=0.5\textwidth]{Use_Case_Diagrams/UC29.png}
                    \caption{UC29 - Inserimento criterio di ricerca ”Denominazione amministrazione AOO”}
            \end{figure}

            \begin{usecase}
                \code{\mdparamIPAOOO}
                \titleUC{Inserimento criterio di ricerca "Codice IPA AOO"}
                \actor{Utente}
                \pre{
                    \item L'utente ha caricato un DIP
                    \item L'utente sta visualizzando una lista di documenti contenuti nel DIP
                    \item L'utente ha configurato il criterio di ricerca "Ruolo" con un valore tra
                    \begin{itemize}
                        \item "Altro"
                        \item "Autore"
                        \item "Destinatario"
                        \item "Mittente"
                    \end{itemize}
                    \item L'utente ha configurato il criterio di ricerca "Tipo soggetto" con il valore "PAI (Amministrazione pubblica italiana)"
                }
                \post{
                    \item L'utente ha configurato un criterio di ricerca
                }
                \main{
                    \item L'utente seleziona la funzionalità di aggiunta criterio di ricerca
                    \item L'utente seleziona il criterio "Codice IPA AOO"
                    \item L'utente inserisce il valore
                }
            \end{usecase}
            \begin{figure}[H]
                    \centering
                    \includegraphics[width=0.5\textwidth]{Use_Case_Diagrams/UC30.png}
                    \caption{UC30 - Inserimento criterio di ricerca ”Codice IPA AOO”}
            \end{figure}

            \begin{usecase}
                \code{\mdparamDenAmmUOR}
                \titleUC{Inserimento criterio di ricerca "Denominazione amministrazione UOR"}
                \actor{Utente}
                \pre{
                    \item L'utente ha caricato un DIP
                    \item L'utente sta visualizzando una lista di documenti contenuti nel DIP
                    \item L'utente ha configurato il criterio di ricerca "Ruolo" con un valore tra
                    \begin{itemize}
                        \item "Altro"
                        \item "Autore"
                        \item "Destinatario"
                        \item "Mittente"
                    \end{itemize}
                    \item L'utente ha configurato il criterio di ricerca "Tipo soggetto" con il valore "PAI (Amministrazione pubblica italiana)"
                }
                \post{
                    \item L'utente ha configurato un criterio di ricerca
                }
                \main{
                    \item L'utente seleziona la funzionalità di aggiunta criterio di ricerca
                    \item L'utente seleziona il criterio "Denominazione amministrazione UOR"
                    \item L'utente inserisce il valore
                }
            \end{usecase}
            \begin{figure}[H]
                    \centering
                    \includegraphics[width=0.5\textwidth]{Use_Case_Diagrams/UC31.png}
                    \caption{UC31 - Inserimento criterio di ricerca ”Denominazione amministrazione UOR”}
            \end{figure}

            \begin{usecase}
                \code{\mdparamIPAUOR}
                \titleUC{Inserimento criterio di ricerca "Codice IPA UOR"}
                \actor{Utente}
                \pre{
                    \item L'utente ha caricato un DIP
                    \item L'utente sta visualizzando una lista di documenti contenuti nel DIP
                    \item L'utente ha configurato il criterio di ricerca "Ruolo" con un valore tra
                    \begin{itemize}
                        \item "Altro"
                        \item "Autore"
                        \item "Destinatario"
                        \item "Mittente"
                    \end{itemize}
                    \item L'utente ha configurato il criterio di ricerca "Tipo soggetto" con il valore "PAI (Amministrazione pubblica italiana)"
                }
                \post{
                    \item L'utente ha configurato un criterio di ricerca
                }
                \main{
                    \item L'utente seleziona la funzionalità di aggiunta criterio di ricerca
                    \item L'utente seleziona il criterio "Codice IPA UOR"
                    \item L'utente inserisce il valore
                }
            \end{usecase}
            \begin{figure}[H]
                    \centering
                    \includegraphics[width=0.5\textwidth]{Use_Case_Diagrams/UC32.png}
                    \caption{UC32 - Inserimento criterio di ricerca ”Codice IPA UOR”}
            \end{figure}

            \begin{usecase}
                \code{\mdparamDenSis}
                \titleUC{Inserimento criterio di ricerca "Denominazione sistema"}
                \actor{Utente}
                \pre{
                    \item L'utente ha caricato un DIP
                    \item L'utente sta visualizzando una lista di documenti contenuti nel DIP
                    \item L'utente ha configurato il criterio di ricerca "Ruolo" con valore "Produttore"
                    \item L'utente ha configurato il criterio di ricerca "Tipo soggetto" con il valore "SW (Documenti prodotti automaticamente)"
                }
                \post{
                    \item L'utente ha configurato un criterio di ricerca
                }
                \main{
                    \item L'utente seleziona la funzionalità di aggiunta criterio di ricerca
                    \item L'utente seleziona il criterio "Denominazione sistema"
                    \item L'utente inserisce il valore
                }
            \end{usecase}
            \begin{figure}[H]
                    \centering
                    \includegraphics[width=0.5\textwidth]{Use_Case_Diagrams/UC33.png}
                    \caption{UC33 - Inserimento criterio di ricerca ”Denominazione sistema”}
            \end{figure}
            
            \begin{usecase}
                \code{\mdparamOgg}
                \titleUC{Inserimento criterio di ricerca "Oggetto"}
                \actor{Utente}
                \pre{
                    \item L'utente ha caricato un DIP
                    \item L'utente sta visualizzando una lista di documenti contenuti nel DIP
                }
                \post{
                    \item L'utente ha configurato un criterio di ricerca
                }
                \main{
                    \item L'utente seleziona la funzionalità di aggiunta criterio di ricerca
                    \item L'utente seleziona il criterio "Oggetto"
                    \item L'utente inserisce il valore
                }
            \end{usecase}
            \begin{figure}[H]
                    \centering
                    \includegraphics[width=0.5\textwidth]{Use_Case_Diagrams/UC34.png}
                    \caption{UC34 - Inserimento criterio di ricerca ”Oggetto”}
            \end{figure}

            \begin{usecase}
                \code{\mdparamParChi}
                \titleUC{Inserimento criterio di ricerca "Parola chiave"}
                \actor{Utente}
                \pre{
                    \item L'utente ha caricato un DIP
                    \item L'utente sta visualizzando una lista di documenti contenuti nel DIP
                }
                \post{
                    \item L'utente ha configurato un criterio di ricerca
                }
                \main{
                    \item L'utente seleziona la funzionalità di aggiunta criterio di ricerca
                    \item L'utente seleziona il criterio "Parola chiave"
                    \item L'utente inserisce il valore
                }
            \end{usecase}
            \begin{figure}[H]
                    \centering
                    \includegraphics[width=0.5\textwidth]{Use_Case_Diagrams/UC35.png}
                    \caption{UC35 - Inserimento criterio di ricerca ”Parola chiave”}
            \end{figure}

            \begin{usecase}
                \code{\mdparamNumAll}
                \titleUC{Inserimento criterio di ricerca "Numero allegati"}
                \actor{Utente}
                \pre{
                    \item L'utente ha caricato un DIP
                    \item L'utente sta visualizzando una lista di documenti contenuti nel DIP
                }
                \post{
                    \item L'utente ha configurato un criterio di ricerca
                }
                \main{
                    \item L'utente seleziona la funzionalità di aggiunta criterio di ricerca
                    \item L'utente seleziona il criterio "Numero allegati"
                    \item L'utente inserisce il valore
                }
            \end{usecase}
            \begin{figure}[H]
                    \centering
                    \includegraphics[width=0.5\textwidth]{Use_Case_Diagrams/UC36.png}
                    \caption{UC36 - Inserimento criterio di ricerca ”Numero allegati”}
            \end{figure}

            \begin{usecase}
                \code{\mdparamImpCritAll}
                \titleUC{Inserimento criterio di ricerca "Impronta crittografica allegato"}
                \actor{Utente}
                \pre{
                    \item L'utente ha caricato un DIP
                    \item L'utente sta visualizzando una lista di documenti contenuti nel DIP
                }
                \post{
                    \item L'utente ha configurato un criterio di ricerca
                }
                \main{
                    \item L'utente seleziona la funzionalità di aggiunta criterio di ricerca
                    \item L'utente seleziona il criterio "Impronta crittografica allegato"
                    \item L'utente inserisce il digest del documento (\textbf{UC\mdparamImpCrit.1})
                    \item L'utente seleziona l'algoritmo applicato (\textbf{UC\mdparamImpCrit.2})
                }
            \end{usecase}

            \begin{usecase}
                \code{\mdparamIdDocAll}
                \titleUC{Inserimento criterio di ricerca "Identificativo allegato"}
                \actor{Utente}
                \pre{
                    \item L'utente ha caricato un DIP
                    \item L'utente sta visualizzando una lista di documenti contenuti nel DIP
                }
                \post{
                    \item L'utente ha configurato un criterio di ricerca
                }
                \main{
                    \item L'utente seleziona la funzionalità di aggiunta criterio di ricerca
                    \item L'utente seleziona il criterio "Identificativo allegato"
                    \item L'utente inserisce il valore
                }
            \end{usecase}

            \begin{usecase}
                \code{\mdparamDesAll}
                \titleUC{Inserimento criterio di ricerca "Descrizione allegato"}
                \actor{Utente}
                \pre{
                    \item L'utente ha caricato un DIP
                    \item L'utente sta visualizzando una lista di documenti contenuti nel DIP
                }
                \post{
                    \item L'utente ha configurato un criterio di ricerca
                }
                \main{
                    \item L'utente seleziona la funzionalità di aggiunta criterio di ricerca
                    \item L'utente seleziona il criterio "Descrizione allegato"
                    \item L'utente inserisce il valore
                }
            \end{usecase}
            \begin{figure}[H]
                    \centering
                    \includegraphics[width=0.5\textwidth]{Use_Case_Diagrams/UC39.png}
                    \caption{UC39 - Inserimento criterio di ricerca ”Descrizione allegato”}
            \end{figure}

            \begin{usecase}
                \code{\mdparamIndCla}
                \titleUC{Inserimento criterio di ricerca "Indice di classificazione"}
                \actor{Utente}
                \pre{
                    \item L'utente ha caricato un DIP
                    \item L'utente sta visualizzando una lista di documenti contenuti nel DIP
                }
                \post{
                    \item L'utente ha configurato un criterio di ricerca
                }
                \main{
                    \item L'utente seleziona la funzionalità di aggiunta criterio di ricerca
                    \item L'utente seleziona il criterio "Indice di classificazione"
                    \item L'utente inserisce il valore
                }
            \end{usecase}
            \begin{figure}[H]
                    \centering
                    \includegraphics[width=0.5\textwidth]{Use_Case_Diagrams/UC40.png}
                    \caption{UC40 - Inserimento criterio di ricerca ”Indice di classificazione”}
            \end{figure}

            \begin{usecase}
                \code{\mdparamDesCla}
                \titleUC{Inserimento criterio di ricerca "Descrizione dell'indice di classificazione"}
                \actor{Utente}
                \pre{
                    \item L'utente ha caricato un DIP
                    \item L'utente sta visualizzando una lista di documenti contenuti nel DIP
                }
                \post{
                    \item L'utente ha configurato un criterio di ricerca
                }
                \main{
                    \item L'utente seleziona la funzionalità di aggiunta criterio di ricerca
                    \item L'utente seleziona il criterio "Descrizione dell'indice di classificazione"
                    \item L'utente inserisce il valore
                }
            \end{usecase}
            \begin{figure}[H]
                    \centering
                    \includegraphics[width=0.5\textwidth]{Use_Case_Diagrams/UC41.png}
                    \caption{UC41 - Inserimento criterio di ricerca ”Descrizione dell’indice di classificazione”}
            \end{figure}

            \begin{usecase}
                \code{\mdparamPiaCla}
                \titleUC{Inserimento criterio di ricerca "Piano di classificazione"}
                \actor{Utente}
                \pre{
                    \item L'utente ha caricato un DIP
                    \item L'utente sta visualizzando una lista di documenti contenuti nel DIP
                }
                \post{
                    \item L'utente ha configurato un criterio di ricerca
                }
                \main{
                    \item L'utente seleziona la funzionalità di aggiunta criterio di ricerca
                    \item L'utente seleziona il criterio "Piano di classificazione"
                    \item L'utente inserisce il valore
                }
            \end{usecase}
            \begin{figure}[H]
                    \centering
                    \includegraphics[width=0.5\textwidth]{Use_Case_Diagrams/UC42.png}
                    \caption{UC42 - Inserimento criterio di ricerca ”Piano di classificazione”}
            \end{figure}

            \begin{usecase}
                \code{\mdparamRis}
                \titleUC{Inserimento criterio di ricerca "Riservato"}
                \actor{Utente}
                \pre{
                    \item L'utente ha caricato un DIP
                    \item L'utente sta visualizzando una lista di documenti contenuti nel DIP
                }
                \post{
                    \item L'utente ha configurato un criterio di ricerca
                }
                \main{
                    \item L'utente seleziona la funzionalità di aggiunta criterio di ricerca
                    \item L'utente seleziona il criterio "Riservato"
                    \item L'utente seleziona il valore tra le opzioni
                    \begin{itemize}
                        \item "Vero"
                        \item "Falso"
                    \end{itemize}
                }
            \end{usecase}
            \begin{figure}[H]
                    \centering
                    \includegraphics[width=0.5\textwidth]{Use_Case_Diagrams/UC43.png}
                    \caption{UC43 - Inserimento criterio di ricerca ”Riservato”}
            \end{figure}

            \begin{usecase}
                \code{\mdparamFor}
                \titleUC{Inserimento criterio di ricerca "Formato"}
                \actor{Utente}
                \pre{
                    \item L'utente ha caricato un DIP
                    \item L'utente sta visualizzando una lista di documenti contenuti nel DIP
                }
                \post{
                    \item L'utente ha configurato un criterio di ricerca
                }
                \main{
                    \item L'utente seleziona la funzionalità di aggiunta criterio di ricerca
                    \item L'utente seleziona il criterio "Formato"
                    \item L'utente seleziona il valore tra le opzioni disponibili
                }
            \end{usecase}
            \begin{figure}[H]
                    \centering
                    \includegraphics[width=0.5\textwidth]{Use_Case_Diagrams/UC44.png}
                    \caption{UC44 - Inserimento criterio di ricerca ”Formato”}
            \end{figure}

            \begin{usecase}
                \code{\mdparamNomPro}
                \titleUC{Inserimento criterio di ricerca "Nome prodotto software"}
                \actor{Utente}
                \pre{
                    \item L'utente ha caricato un DIP
                    \item L'utente sta visualizzando una lista di documenti contenuti nel DIP
                }
                \post{
                    \item L'utente ha configurato un criterio di ricerca
                }
                \main{
                    \item L'utente seleziona la funzionalità di aggiunta criterio di ricerca
                    \item L'utente seleziona il criterio "Nome prodotto software"
                    \item L'utente inserisce il valore
                }
            \end{usecase}
            \begin{figure}[H]
                    \centering
                    \includegraphics[width=0.5\textwidth]{Use_Case_Diagrams/UC45.png}
                    \caption{UC45 - Inserimento criterio di ricerca ”Nome prodotto software”}
            \end{figure}

            \begin{usecase}
                \code{\mdparamVerPro}
                \titleUC{Inserimento criterio di ricerca "Versione prodotto software"}
                \actor{Utente}
                \pre{
                    \item L'utente ha caricato un DIP
                    \item L'utente sta visualizzando una lista di documenti contenuti nel DIP
                }
                \post{
                    \item L'utente ha configurato un criterio di ricerca
                }
                \main{
                    \item L'utente seleziona la funzionalità di aggiunta criterio di ricerca
                    \item L'utente seleziona il criterio "Versione prodotto software"
                    \item L'utente inserisce il valore
                }
            \end{usecase}
            \begin{figure}[H]
                    \centering
                    \includegraphics[width=0.5\textwidth]{Use_Case_Diagrams/UC46.png}
                    \caption{UC46 - Inserimento criterio di ricerca ”Versione prodotto software”}
            \end{figure}

            \begin{usecase}
                \code{\mdparamProd}
                \titleUC{Inserimento criterio di ricerca "Produttore software"}
                \actor{Utente}
                \pre{
                    \item L'utente ha caricato un DIP
                    \item L'utente sta visualizzando una lista di documenti contenuti nel DIP
                }
                \post{
                    \item L'utente ha configurato un criterio di ricerca
                }
                \main{
                    \item L'utente seleziona la funzionalità di aggiunta criterio di ricerca
                    \item L'utente seleziona il criterio "Produttore software"
                    \item L'utente inserisce il valore
                }
            \end{usecase}
            \begin{figure}[H]
                    \centering
                    \includegraphics[width=0.5\textwidth]{Use_Case_Diagrams/UC47.png}
                    \caption{UC47 - Inserimento criterio di ricerca ”Produttore software”}
            \end{figure}

            \begin{usecase}
                \code{\mdparamFirDig}
                \titleUC{Inserimento criterio di ricerca "Firmato digitalmente"}
                \actor{Utente}
                \pre{
                    \item L'utente ha caricato un DIP
                    \item L'utente sta visualizzando una lista di documenti contenuti nel DIP
                }
                \post{
                    \item L'utente ha configurato un criterio di ricerca
                }
                \main{
                    \item L'utente seleziona la funzionalità di aggiunta criterio di ricerca
                    \item L'utente seleziona il criterio "Firmato digitalmente"
                    \item L'utente seleziona il valore tra le opzioni
                    \begin{itemize}
                        \item "Vero"
                        \item "Falso"
                    \end{itemize}
                }
            \end{usecase}
            \begin{figure}[H]
                    \centering
                    \includegraphics[width=0.5\textwidth]{Use_Case_Diagrams/UC48.png}
                    \caption{UC48 - Inserimento criterio di ricerca ”Firmato digitalmente”}
            \end{figure}

            \begin{usecase}
                \code{\mdparamSigEle}
                \titleUC{Inserimento criterio di ricerca "Sigillato elettronicamente"}
                \actor{Utente}
                \pre{
                    \item L'utente ha caricato un DIP
                    \item L'utente sta visualizzando una lista di documenti contenuti nel DIP
                }
                \post{
                    \item L'utente ha configurato un criterio di ricerca
                }
                \main{
                    \item L'utente seleziona la funzionalità di aggiunta criterio di ricerca
                    \item L'utente seleziona il criterio "Sigillato elettronicamente"
                    \item L'utente seleziona il valore tra le opzioni
                    \begin{itemize}
                        \item "Vero"
                        \item "Falso"
                    \end{itemize}
                }
            \end{usecase}
            \begin{figure}[H]
                    \centering
                    \includegraphics[width=0.5\textwidth]{Use_Case_Diagrams/UC49.png}
                    \caption{UC49 - Inserimento criterio di ricerca ”Sigillato elettronicamente”}
            \end{figure}

            \begin{usecase}
                \code{\mdparamMarTem}
                \titleUC{Inserimento criterio di ricerca "Marcatura temporale"}
                \actor{Utente}
                \pre{
                    \item L'utente ha caricato un DIP
                    \item L'utente sta visualizzando una lista di documenti contenuti nel DIP
                }
                \post{
                    \item L'utente ha configurato un criterio di ricerca
                }
                \main{
                    \item L'utente seleziona la funzionalità di aggiunta criterio di ricerca
                    \item L'utente seleziona il criterio "Marcatura temporale"
                    \item L'utente seleziona il valore tra le opzioni
                    \begin{itemize}
                        \item "Vero"
                        \item "Falso"
                    \end{itemize}
                }
            \end{usecase}
            \begin{figure}[H]
                    \centering
                    \includegraphics[width=0.5\textwidth]{Use_Case_Diagrams/UC50.png}
                    \caption{UC50 - Inserimento criterio di ricerca ”Marcatura temporale”}
            \end{figure}

            \begin{usecase}
                \code{\mdparamConCop}
                \titleUC{Inserimento criterio di ricerca "Conformità copie immagine su supporto informatico"}
                \actor{Utente}
                \pre{
                    \item L'utente ha caricato un DIP
                    \item L'utente sta visualizzando una lista di documenti contenuti nel DIP
                }
                \post{
                    \item L'utente ha configurato un criterio di ricerca
                }
                \main{
                    \item L'utente seleziona la funzionalità di aggiunta criterio di ricerca
                    \item L'utente seleziona il criterio "Conformità copie immagine su supporto informatico"
                    \item L'utente seleziona il valore tra le opzioni
                    \begin{itemize}
                        \item "Vero"
                        \item "Falso"
                    \end{itemize}
                }
            \end{usecase}
            \begin{figure}[H]
                    \centering
                    \includegraphics[width=0.5\textwidth]{Use_Case_Diagrams/UC51.png}
                    \caption{UC51 - Inserimento criterio di ricerca ”Conformità copie immagine su supporto informatico”}
            \end{figure}

            \begin{usecase}
                \code{\mdparamTipAgg}
                \titleUC{Inserimento criterio di ricerca "Tipo aggregazione"}
                \actor{Utente}
                \pre{
                    \item L'utente ha caricato un DIP
                    \item L'utente sta visualizzando una lista di documenti contenuti nel DIP
                }
                \post{
                    \item L'utente ha configurato un criterio di ricerca
                }
                \main{
                    \item L'utente seleziona la funzionalità di aggiunta criterio di ricerca
                    \item L'utente seleziona il criterio "Tipo aggregazione"
                    \item L'utente seleziona il valore tra le opzioni
                    \begin{itemize}
                        \item "Fascicolo"
                        \item "Serie documentale"
                        \item "Serie di fascicoli"
                    \end{itemize}
                }
            \end{usecase}
            \begin{figure}[H]
                    \centering
                    \includegraphics[width=0.5\textwidth]{Use_Case_Diagrams/UC52.png}
                    \caption{UC52 - Inserimento criterio di ricerca ”Tipo aggregazione”}
            \end{figure}

            \begin{usecase}
                \code{\mdparamIdAgg}
                \titleUC{Inserimento criterio di ricerca "Identificativo aggregazione"}
                \actor{Utente}
                \pre{
                    \item L'utente ha caricato un DIP
                    \item L'utente sta visualizzando una lista di documenti contenuti nel DIP
                }
                \post{
                    \item L'utente ha configurato un criterio di ricerca
                }
                \main{
                    \item L'utente seleziona la funzionalità di aggiunta criterio di ricerca
                    \item L'utente seleziona il criterio "Identificativo aggregazione"
                    \item L'utente inserisce il valore
                }
            \end{usecase}
            \begin{figure}[H]
                    \centering
                    \includegraphics[width=0.5\textwidth]{Use_Case_Diagrams/UC53.png}
                    \caption{UC53 - Inserimento criterio di ricerca ”Identificativo aggregazione”}
            \end{figure}

            \begin{usecase}
                \code{\mdparamTipFas}
                \titleUC{Inserimento criterio di ricerca "Tipologia fascicolo"}
                \actor{Utente}
                \pre{
                    \item L'utente ha caricato un DIP
                    \item L'utente sta visualizzando una lista di documenti contenuti nel DIP
                    \item L'utente ha configurato il criterio di ricerca "Tipo aggregazione" con valore "Fascicolo"
                }
                \post{
                    \item L'utente ha configurato un criterio di ricerca
                }
                \main{
                    \item L'utente seleziona la funzionalità di aggiunta criterio di ricerca
                    \item L'utente seleziona il criterio "Tipologia fascicolo"
                    \item L'utente seleziona il valore tra le opzioni
                    \begin{itemize}
                        \item "Affare"
                        \item "Attività"
                        \item "Persona fisica"
                        \item "Persona giuridica"
                        \item "Procedimento amministrativo"
                    \end{itemize}
                }
            \end{usecase}
            \begin{figure}[H]
                    \centering
                    \includegraphics[width=0.5\textwidth]{Use_Case_Diagrams/UC54.png}
                    \caption{UC54 - Inserimento criterio di ricerca ”Tipologia fascicolo”}
            \end{figure}

            \begin{usecase}
                \code{\mdparamNomDoc}
                \titleUC{Inserimento criterio di ricerca "Nome del documento"}
                \actor{Utente}
                \pre{
                    \item L'utente ha caricato un DIP
                    \item L'utente sta visualizzando una lista di documenti contenuti nel DIP
                }
                \post{
                    \item L'utente ha configurato un criterio di ricerca
                }
                \main{
                    \item L'utente seleziona la funzionalità di aggiunta criterio di ricerca
                    \item L'utente seleziona il criterio "Nome del documento"
                    \item L'utente inserisce il valore
                }
            \end{usecase}
            \begin{figure}[H]
                    \centering
                    \includegraphics[width=0.5\textwidth]{Use_Case_Diagrams/UC55.png}
                    \caption{UC55 - Inserimento criterio di ricerca ”Nome del documento”}
            \end{figure}

            \begin{usecase}
                \code{\mdparamVerDoc}
                \titleUC{Inserimento criterio di ricerca "Versione del documento"}
                \actor{Utente}
                \pre{
                    \item L'utente ha caricato un DIP
                    \item L'utente sta visualizzando una lista di documenti contenuti nel DIP
                }
                \post{
                    \item L'utente ha configurato un criterio di ricerca
                }
                \main{
                    \item L'utente seleziona la funzionalità di aggiunta criterio di ricerca
                    \item L'utente seleziona il criterio "Versione del documento"
                    \item L'utente inserisce il valore
                }
            \end{usecase}
            \begin{figure}[H]
                    \centering
                    \includegraphics[width=0.5\textwidth]{Use_Case_Diagrams/UC56.png}
                    \caption{UC56 - Inserimento criterio di ricerca ”Versione del documento”}
            \end{figure}

             \begin{usecase}
                \code{\mdparamTipMod}
                \titleUC{Inserimento criterio di ricerca "Tipo modifica"}
                \actor{Utente}
                \pre{
                    \item L'utente ha caricato un DIP
                    \item L'utente sta visualizzando una lista di documenti contenuti nel DIP
                }
                \post{
                    \item L'utente ha configurato un criterio di ricerca
                }
                \main{
                    \item L'utente seleziona la funzionalità di aggiunta criterio di ricerca
                    \item L'utente seleziona il criterio "Tipo modifica"
                    \item L'utente seleziona il valore tra le opzioni
                    \begin{itemize}
                        \item "Annulamento"
                        \item "Annotazione"
                        \item "Integrazione"
                        \item "Rettifica"
                    \end{itemize}
                }
            \end{usecase}
            \begin{figure}[H]
                    \centering
                    \includegraphics[width=0.5\textwidth]{Use_Case_Diagrams/UC57.png}
                    \caption{UC57 - Inserimento criterio di ricerca ”Tipo modifica”}
            \end{figure}

            \begin{usecase}
                \code{\mdparamDatMod}
                \titleUC{Inserimento criterio di ricerca "Data modifica"}
                \actor{Utente}
                \pre{
                    \item L'utente ha caricato un DIP
                    \item L'utente sta visualizzando una lista di documenti contenuti nel DIP
                }
                \post{
                    \item L'utente ha configurato un criterio di ricerca
                }
                \main{
                    \item L'utente seleziona la funzionalità di aggiunta criterio di ricerca
                    \item L'utente seleziona il criterio "Data modifica"
                    \item L'utente inserisce il valore
                }
            \end{usecase}
            \begin{figure}[H]
                    \centering
                    \includegraphics[width=0.5\textwidth]{Use_Case_Diagrams/UC58.png}
                    \caption{UC58 - Inserimento criterio di ricerca ”Data modifica”}
            \end{figure}

            \begin{usecase}
                \code{\mdparamOraMod}
                \titleUC{Inserimento criterio di ricerca "Ora modifica"}
                \actor{Utente}
                \pre{
                    \item L'utente ha caricato un DIP
                    \item L'utente sta visualizzando una lista di documenti contenuti nel DIP
                }
                \post{
                    \item L'utente ha configurato un criterio di ricerca
                }
                \main{
                    \item L'utente seleziona la funzionalità di aggiunta criterio di ricerca
                    \item L'utente seleziona il criterio "Ora modifica"
                    \item L'utente inserisce il valore
                }
            \end{usecase}
            \begin{figure}[H]
                    \centering
                    \includegraphics[width=0.5\textwidth]{Use_Case_Diagrams/UC59.png}
                    \caption{UC59 - Inserimento criterio di ricerca ”Ora modifica”}
            \end{figure}

            \begin{usecase}
                \code{\mdparamTemCon}
                \titleUC{Inserimento criterio di ricerca "Tempo di conservazione"}
                \actor{Utente}
                \pre{
                    \item L'utente ha caricato un DIP
                    \item L'utente sta visualizzando una lista di documenti contenuti nel DIP
                }
                \post{
                    \item L'utente ha configurato un criterio di ricerca
                }
                \main{
                    \item L'utente seleziona la funzionalità di aggiunta criterio di ricerca
                    \item L'utente seleziona il criterio "Tempo di conservazione"
                    \item L'utente inserisce il valore
                }
            \end{usecase}
            \begin{figure}[H]
                    \centering
                    \includegraphics[width=0.5\textwidth]{Use_Case_Diagrams/UC60.png}
                    \caption{UC60 - Inserimento criterio di ricerca ”Tempo di conservazione”}
            \end{figure}

            \begin{usecase}
                \code{\mdparamNot}
                \titleUC{Inserimento criterio di ricerca "Note"}
                \actor{Utente}
                \pre{
                    \item L'utente ha caricato un DIP
                    \item L'utente sta visualizzando una lista di documenti contenuti nel DIP
                }
                \post{
                    \item L'utente ha configurato un criterio di ricerca
                }
                \main{
                    \item L'utente seleziona la funzionalità di aggiunta criterio di ricerca
                    \item L'utente seleziona il criterio "Note"
                    \item L'utente inserisce il valore
                }
            \end{usecase}
            \begin{figure}[H]
                    \centering
                    \includegraphics[width=0.5\textwidth]{Use_Case_Diagrams/UC61.png}
                    \caption{UC61 - Inserimento criterio di ricerca ”Note”}
            \end{figure}

            % Ricerca tramite metadati
            \begin{usecase}
                \code{\mdsearch}
                \titleUC{Ricerca tramite metadati}
                \actor{Utente}
                \pre{
                    \item L'utente ha caricato un DIP
                    \item L'utente sta visualizzando una lista di documenti contenuti nel DIP
                    \item L'utente ha configurato uno o più criteri di ricerca tra quelli disponibili
                }
                \post{
                    \item L'utente ha avviato una ricerca
                }
                \main{
                    \item L'utente avvia la ricerca
                }
            \end{usecase}
            \begin{figure}[H]
                    \centering
                    \includegraphics[width=0.5\textwidth]{Use_Case_Diagrams/UC62.png}
                    \caption{UC62 - Ricerca tramite metadati}
            \end{figure}

            % Visualizzazione risultati ricerca
            \begin{usecase}
                \code{\searchres}
                \titleUC{Visualizzazione risultati ricerca}
                \actor{Utente}
                \pre{
                    \item L'utente ha caricato un DIP
                    \item L'utente ha configurato un insieme di criteri di ricerca 
                    \item L'utente ha eseguito una ricerca
                }
                \post{
                    \item L'utente vede la lista dei documenti che corrispondono ai criteri di ricerca
                }
                \main{
                    \item L'utente visualizza i documenti corrispondenti ai criteri di ricerca
                }
                \alt{
                    \item[1a.] Nessun documento corrisponde ai criteri di ricerca (\textbf{UC\nosearchres})
                }
                \inc{\ucref{uc:\searchres.1}}
            \end{usecase}
            \begin{figure}[H]
                    \centering
                    \includegraphics[width=0.5\textwidth]{Use_Case_Diagrams/UC63.png}
                    \caption{UC63 - Visualizzazione risultati ricerca}
            \end{figure}
            \begin{usecase}
                \code{\searchres.1}
                \titleUC{Visualizzazione singolo risultato ricerca}
                \actor{Utente}
                \pre{
                    \item L'utente sta visualizzando i risultati della ricerca
                }
                \post{
                    \item L'utente vede e informazioni principali di un documento che corrisponde ai criteri di ricerca
                }
                \main{
                    \item L'utente visualizza le informazioni principali di un documento che corrisponde ai criteri di ricerca, quali
                    \begin{itemize}
                        \item Il nome del file
                        \item Il suo formato
                        \item Identificativo documentale
                        \item Data in cui il documento è stato registrato
                    \end{itemize}
                }
            \end{usecase}

            % Risultato ricerca vuoto
            \begin{usecase}
                \code{\nosearchres}
                \titleUC{Risultato ricerca vuoto}
                \actor{Utente}
                \pre{
                    \item L'utente ha immesso dei criteri di ricerca che non hanno prodotto risultati
                }
                \post{
                    \item L'utente vede un messaggio di avviso
                }
                \main{
                    \item L'utente visualizza un messaggio informativo sull'assenza di risultati
                }
            \end{usecase}
            \begin{figure}[H]
                    \centering
                    \includegraphics[width=0.5\textwidth]{Use_Case_Diagrams/UC64.png}
                    \caption{UC64 - Risultato ricerca vuoto}
            \end{figure}

            % Anteprima documento
            \begin{usecase}
                \code{\docpreview}
                \titleUC{Visualizzazione anteprima documento}
                \actor{Utente}
                \pre{
                    \item L'utente ha caricato un DIP
                    \item L'utente sta visualizzando una lista di documenti contenuti nel DIP
                    \item L'utente ha selezionato un documento dalla lista
                }
                \post{
                    \item L'utente vede l'anteprima del documento selezionato
                }
                \main{
                    \item L'utente seleziona la funzionalità di anteprima del documento selezionato
                    \item L'utente visualizza l'anteprima del documento selezionato
                }
                \alt{
                    \item[2a.] L'utente seleziona un documento con un formato per cui l'anteprima non è supportata (\textbf{UC\nodocpreview})
                }
                \ext{\ucref{uc:\nodocpreview}}
            \end{usecase}
            \begin{figure}[H]
                    \centering
                    \includegraphics[width=0.5\textwidth]{Use_Case_Diagrams/UC65.png}
                    \caption{UC65 - Visualizzazione anteprima documento}
            \end{figure}

            % Anteprima documento non supportata
            \begin{usecase}
                \code{\nodocpreview}
                \titleUC{Anteprima documento non supportata}
                \actor{Utente}
                \pre{
                    \item L'utente ha selezionato la funzione di anteprima per un documento con formato non supportato
                }
                \post{
                    \item L'utente non vede nessuna anteprima
                    \item L'utente vede un messaggio di avviso
                }
                \main{
                    \item L'utente visualizza un messaggio di avviso sull'indisponibilità dell'anteprima
                }
            \end{usecase}

            % Esportazione documento
            \begin{usecase}
                \code{\exportdoc}
                \titleUC{Esportazione documento}
                \actor{Utente}
                \pre{
                    \item L'utente ha caricato un DIP
                    \item L'utente sta visualizzando una lista di documenti contenuti nel DIP
                    \item L'utente ha selezionato un documento dalla lista
                }
                \post{
                    \item Il documento selezionato è salvato sul dispositivo dell'utente
                    \item L'utente vede un messaggio di conferma
                }
                \main{
                    \item L'utente seleziona la funzionalità di esportazione del documento selezionato
                    \item L'utente inserisce il percorso dove salvare il documento (\textbf{UC\exportdoc.1})
                    \item L'utente visualizza un messaggio conferma dell'avvenuto salvataggio
                }
                \alt{
                    \item[2a.] L'utente cancella l'operazione
                    \item[3a.] L'operazione non ha successo (\textbf{UC\exportdocfail})
                }
                \inc{\ucref{uc:\exportdoc.1}}
                \ext{\ucref{uc:\exportdocfail}}
            \end{usecase}
            \begin{figure}[H]
                    \centering
                    \includegraphics[width=0.5\textwidth]{Use_Case_Diagrams/UC67.png}
                    \caption{UC67 - Esportazione documento}
            \end{figure}
            \begin{usecase}
                \code{\exportdoc.1}
                \titleUC{Inserimento percorso di esportazione}
                \actor{Utente}
                \pre{
                    \item L'utente ha selezionato la funzionalità di esportazione di un documento
                }
                \post{
                    \item L'utente ha inserito il percorso per il salvataggio del file
                }
                \main{
                    \item Il sistema richiede il percorso per il salvataggio del file
                    \item L'utente inserisce il percorso per il salvataggio del file
                }
            \end{usecase}

            % Fallimento esportazione documento
            \begin{usecase}
                \code{\exportdocfail}
                \titleUC{Fallimento esportazione documento}
                \actor{Utente}
                \pre{
                    \item L'operazione di salvataggio non ha avuto successo
                }
                \post{
                    \item L'utente vede un messaggio di errore
                }
                \main{
                    \item L'utente visualizza un messaggio di errore riguardo il fallimento dell'esportazione
                }
            \end{usecase}

            % Parametro di ricerca semantica
            \begin{usecase}
                \code{\semparam}
                \titleUC{Configurazione ricerca semantica}
                \actor{Utente}
                \pre{
                    \item L'utente ha caricato un DIP
                    \item L'utente sta visualizzando una lista di documenti contenuti nel DIP
                }
                \post{
                    \item L'utente ha configurato la ricerca semantica
                }
                \main{
                    \item L'utente seleziona la configurazione della ricerca semantica
                    \item L'utente inserisce la query in linguaggio naturale (\textbf{UC\semparam.1})
                }
                \inc{\ucref{uc:\semparam.1}}
            \end{usecase}
            \begin{figure}[H]
                    \centering
                    \includegraphics[width=0.5\textwidth]{Use_Case_Diagrams/UC69.png}
                    \caption{UC69 - Configurazione ricerca semantica}
            \end{figure}
            
            \begin{usecase}
                \code{\semparam.1}
                \titleUC{Inserimento query in linguaggio naturale}
                \actor{Utente}
                \pre{
                    \item L'utente sta configurando la ricerca semantica
                }
                \post{
                    \item L'utente ha inserito la query in linguaggio naturale
                }
                \main{
                    \item L'utente inserisce la query in linguaggio naturale
                }
            \end{usecase}

            % Ricerca semantica
            \begin{usecase}
                \code{\semsearch}
                \titleUC{Ricerca semantica}
                \actor{Utente}
                \secactor{Motore semantico}
                \pre{
                    \item L'utente ha caricato un DIP
                    \item L'utente sta visualizzando una lista di documenti contenuti nel DIP
                    \item L'utente ha configurato la ricerca semantica
                }
                \post{
                    \item L'utente ha avviato una ricerca
                }
                \main{
                    \item L'utente avvia la ricerca
                    \item Il motore di ricerca semantica elabora la ricerca e produce un risultato
                }
            \end{usecase}
            \begin{figure}[H]
                    \centering
                    \includegraphics[width=0.5\textwidth]{Use_Case_Diagrams/UC70.png}
                    \caption{UC70 - Ricerca semantica}
            \end{figure}

            % Stampa documento
            \begin{usecase}
                \code{\printdoc}
                \titleUC{Stampa documento}
                \actor{Utente}
                \secactor{Gestore di stampa SO}
                \pre{
                    \item L'utente ha caricato un DIP
                    \item L'utente sta visualizzando una lista di documenti contenuti nel DIP
                    \item L'utente ha selezionato un documento dalla lista
                }
                \post{
                    \item Il documento selezionato è inviato alla coda di stampa
                    \item L'utente vede un messaggio di conferma
                }
                \main{
                    \item L'utente seleziona la funzionalità di stampa per il documento selezionato
                    \item Il gestore di stampa riceve il documento
                    \item L'utente visualizza un messaggio relativo al completamento dell'operazione
                }
                \alt{
                    \item[3a.] L'operazione di stampa non va a buon fine (\textbf{UC\printdocfail})
                }
                \ext{\ucref{uc:\printdocfail}}
            \end{usecase}
            \begin{figure}[H]
                    \centering
                    \includegraphics[width=0.5\textwidth]{Use_Case_Diagrams/UC71.png}
                    \caption{UC71 - Stampa documento}
            \end{figure}

            % Fallimento stampa documento
            \begin{usecase}
                \code{\printdocfail}
                \titleUC{Fallimento stampa documento}
                \actor{Utente}
                \secactor{Gestore di stampa SO}
                \pre{
                    \item L'operazione di stampa non ha avuto successo
                }
                \post{
                    \item L'utente vede un messaggio di errore
                }
                \main{
                    \item L'utente visualizza un messaggio di errore riguardo il fallimento della stampa
                }
            \end{usecase}

            % Verifica firma digitale
            \begin{usecase}
                \code{\verdigsign}
                \titleUC{Verifica firma digitale documento}
                \actor{Utente}
                \pre{
                    \item L'utente ha caricato un DIP
                    \item L'utente sta visualizzando una lista di documenti contenuti nel DIP
                    \item L'utente ha selezionato un documento dalla lista
                }
                \post{
                    \item L'utente ha verificato la validità della firma digitale
                }
                \main{
                    \item L'utente seleziona la funzionalità di verifica delle firma digitale del documento selezionato
                    \item L'utente avvia la verifica
                    \item L'utente ha verificato la validità della firma digitale
                }
            \end{usecase}
            \begin{figure}[H]
                    \centering
                    \includegraphics[width=0.5\textwidth]{Use_Case_Diagrams/UC73.png}
                    \caption{UC73 - Verifica firma digitale documento}
            \end{figure}

            % Visualizzazione risultato verifica firma digitale
            \begin{usecase}
                \code{\viewverdigsign}
                \titleUC{Visualizzazione risultato verifica firma digitale documento}
                \actor{Utente}
                \pre{
                    \item L'utente ha verificato la validità della firma digitale
                }
                \post{
                    \item L'utente vede il risultato della verifica della firma digitale 
                }
                \main{
                    \item L'utente visualizza il risultato delle verifica della firma digitale (firma valida/non valida)
                }
            \end{usecase}
            \begin{figure}[H]
                    \centering
                    \includegraphics[width=0.5\textwidth]{Use_Case_Diagrams/UC74.png}
                    \caption{UC74 - Visualizzazione risultato verifica firma digitale documento}
            \end{figure}
        }
    }
    \newpage
    \section{Requisiti} {
        In questa sezione vengono illustrati i requisiti alla base di DIPReader. Essi sono suddivisi in:
        \begin{itemize}
            \item \textbf{Requisiti funzionali:} definiscono i servizi e le funzionalità che il sistema deve fornire.
            \item \textbf{Requisiti di qualità:} descrivono le caratteristiche che il sistema deve soddisfare per accertarne la qualità.
            \item \textbf{Requisiti di vincolo:} impongono delle restrizioni al Sistema, ad esempio riguardo le tecnologie da utilizzare.
        \end{itemize}
        Un ulteriore suddivisione dei requisiti, viene fatta sulla base dell'importanza della funzionalità, caratteristica, o vincolo che descrivono. In particolare possono essere:
        \begin{itemize}
            \item \textbf{Obbligatori:} sono essenziali, il sistema non può funzionare correttamente senza di essi.
            \item \textbf{Desiderabili:} migliorano l’usabilità o l’efficienza del sistema, ma la mancanza di questi non impedisce al sistema di funzionare.
            \item \textbf{Opzionali:} sono caratteristiche aggiuntive e negoziabili che possono essere implementate in futuro, in quanto relativamente utili o per circostanze di limitata disponibilità di tempo e risorse.
        \end{itemize}
        Ogni requisito viene identificato univocamente da un codice nel formato \textbf{RX-Y-Z}:
        \begin{itemize}
            \item \textbf{X} rappresenta il numero del requisito.
            \item \textbf{Y} può assumere i valori \textbf{F}, \textbf{Q}, \textbf{V}, per rappresentare, rispettivamente, un requisito \textbf{Funzionale}, di \textbf{Qualità} o di \textbf{Vincolo}.
            \item \textbf{Z} può assumere i valori \textbf{Ob}, \textbf{De}, \textbf{Op}, per rappresentare, rispettivamente, un requisito \textbf{Obbligatorio}, \textbf{Desiderabile} o \textbf{Opzionale}.
        \end{itemize}
        \subsection{Requisiti Funzionali} {
            \rowcolors{2}{lightgray}{white}
            \begin{longtable}{|l|P{0.5\linewidth}|P{0.3\linewidth}|}
                \caption{Requisiti Funzionali} \\
                
                \hline
                \rowcolor{lightgray}
                \textbf{Codice} & \textbf{Descrizione} & \textbf{Fonti} \\
                \hline
                \endfirsthead
                
                \hline
                \rowcolor{lightgray}
                \textbf{Codice} & \textbf{Descrizione} & \textbf{Fonti} \\
                \hline
                \endhead
                
                \textbf{R1-F-Ob} & L'utente deve poter selezionare un DIP da caricare nel Sistema & 
                \ucref{uc:\loaddip} \newline
                \ucref{uc:\loaddip.1} \\
                \hline
                \textbf{R2-F-Ob} & L'utente deve ricevere un messaggio di errore se il caricamento del DIP non avviene correttamente & 
                \ucref{uc:\remloaddip} \newline
                \ucref{uc:\errloaddip} \\
                \hline
                \textbf{R3-F-Ob} & L'utente deve poter rimuovere un DIP dal Sistema & 
                \ucref{uc:\removedip} \\
                \hline
                \textbf{R4-F-Ob} & L'utente deve poter visualizzare il contenuto e le informazioni del DIP & \ucref{uc:\viewdip} \newline
                \ucref{uc:\viewdip.1} \newline
                \ucref{uc:\viewdip.2} \\
                \hline
                \textbf{R5-F-Ob} & L'utente deve poter visualizzare i singoli documenti, e la lista dei loro allegati, contenuti nel DIP & 
                \ucref{uc:\viewdip} \newline
                \ucref{uc:\viewdip.2.1} \newline
                \ucref{uc:\viewdip.2.2} \newline
                \ucref{uc:\viewdip.2.2.1} \\
                \hline
                \textbf{R6-F-Ob} & L'utente deve poter visualizzare i singoli allegati di un documento & \ucref{uc:\viewdip} \newline
                \ucref{uc:\viewdip.2.2.1} \newline
                \ucref{uc:\viewdip.2.2.1.1} \\
                \hline
                \textbf{R7-F-Ob} & L'utente deve poter visualizzare i dettagli di un documento & \ucref{uc:\viewdetdoc} \\
                \hline
                \textbf{R8-F-Ob} & L'utente deve poter visualizzare le informazioni di conservazione di un DIP & \ucref{uc:\viewconsinfo} \\
                \hline
                \textbf{R9-F-Ob} & Il Sistema deve eseguire i controlli sulla validità e lo stato del processo di conservazione & \ucref{uc:\viewconsinfo} \\
                \hline
                \textbf{R10-F-De} & L'utente deve poter impostare il criterio "Impronta crittografica" per cercare un documento & 
                \ucref{uc:\mdparamImpCrit} \newline
                \ucref{uc:\mdparamImpCrit.1} \newline
                \ucref{uc:\mdparamImpCrit.2} \\
                \hline
                \textbf{R11-F-Ob} & L'utente deve poter impostare il criterio "Identificativo documento" per cercare un documento & 
                \ucref{uc:\mdparamIdDoc} \\
                \hline
                \textbf{R12-F-De} & L'utente deve poter impostare il criterio "Modalità di formazione" per cercare un documento & 
                \ucref{uc:\mdparamModFor} \\
                \hline
                \textbf{R13-F-Ob} & L'utente deve poter impostare il criterio "Tipologia documentale" per cercare un documento & 
                \ucref{uc:\mdparamTipDoc} \\
                \hline
                \textbf{R14-F-De} & L'utente deve poter impostare il criterio "Tipologia di flusso" per cercare un documento & 
                \ucref{uc:\mdparamTipFlu} \\
                \hline
                \textbf{R15-F-De} & L'utente deve poter impostare il criterio "Tipo registro" per cercare un documento & 
                \ucref{uc:\mdparamTipReg} \\
                \hline
                \textbf{R16-F-Ob} & L'utente deve poter impostare il criterio "Data registrazione" per cercare un documento & 
                \ucref{uc:\mdparamDatReg} \\
                \hline
                \textbf{R17-F-Ob} & L'utente deve poter impostare il criterio "Ora registrazione" per cercare un documento & 
                \ucref{uc:\mdparamOraReg} \\
                \hline
                \textbf{R18-F-Ob} & L'utente deve poter impostare il criterio "Numero documento" per cercare un documento & 
                \ucref{uc:\mdparamNumDoc} \\
                \hline
                \textbf{R19-F-Ob} & L'utente deve poter impostare il criterio "Codice registro" per cercare un documento & 
                \ucref{uc:\mdparamCodReg} \\
                \hline
                \textbf{R20-F-Ob} & L'utente deve poter impostare il criterio "Ruolo" per cercare un documento & 
                \ucref{uc:\mdparamRuolo} \\
                \hline
                \textbf{R21-F-Ob} & L'utente deve poter impostare il criterio "Tipo soggetto" per cercare un documento & 
                \ucref{uc:\mdparamTipSog} \\
                \hline
                \textbf{R22-F-Ob} & L'utente deve poter impostare il criterio "Cognome" per cercare un documento & 
                \ucref{uc:\mdparamCog} \\
                \hline
                \textbf{R23-F-Ob} & L'utente deve poter impostare il criterio "Nome" per cercare un documento & 
                \ucref{uc:\mdparamNom} \\
                \hline
                \textbf{R24-F-Ob} & L'utente deve poter impostare il criterio "Codice Fiscale" per cercare un documento & 
                \ucref{uc:\mdparamCF} \\
                \hline
                \textbf{R25-F-De} & L'utente deve poter impostare il criterio "Indirizzo digitale" per cercare un documento & 
                \ucref{uc:\mdparamIndDig} \\
                \hline
                \textbf{R26-F-De} & L'utente deve poter impostare il criterio "Denominazione organizzazione" per cercare un documento & 
                \ucref{uc:\mdparamDenOrg} \\
                \hline
                \textbf{R27-F-De} & L'utente deve poter impostare il criterio "Partita IVA" per cercare un documento & 
                \ucref{uc:\mdparamCFPI} \\
                \hline
                \textbf{R28-F-De} & L'utente deve poter impostare il criterio "Denominazione ufficio" per cercare un documento & 
                \ucref{uc:\mdparamDenUff} \\
                \hline
                \textbf{R29-F-De} & L'utente deve poter impostare il criterio "Denominazione amministrazione" per cercare un documento & 
                \ucref{uc:\mdparamDenAmm} \\
                \hline
                \textbf{R30-F-De} & L'utente deve poter impostare il criterio "Codice IPA" per cercare un documento & 
                \ucref{uc:\mdparamIPA} \\
                \hline
                \textbf{R31-F-De} & L'utente deve poter impostare il criterio "Denominazione amministrazione AOO" per cercare un documento & 
                \ucref{uc:\mdparamDenAmmAOO} \\
                \hline
                \textbf{R32-F-De} & L'utente deve poter impostare il criterio "Codice IPA AOO" per cercare un documento & 
                \ucref{uc:\mdparamIPAOOO} \\
                \hline
                \textbf{R33-F-De} & L'utente deve poter impostare il criterio "Denominazione amministrazione UOR" per cercare un documento & 
                \ucref{uc:\mdparamDenAmmUOR} \\
                \hline
                \textbf{R34-F-De} & L'utente deve poter impostare il criterio "Codice IPA UOR" per cercare un documento & 
                \ucref{uc:\mdparamIPAUOR} \\
                \hline
                \textbf{R35-F-De} & L'utente deve poter impostare il criterio "Denominazione sistema" per cercare un documento & 
                \ucref{uc:\mdparamDenSis} \\
                \hline
                \textbf{R36-F-Ob} & L'utente deve poter impostare il criterio "Oggetto" per cercare un documento & 
                \ucref{uc:\mdparamOgg} \\
                \hline
                \textbf{R37-F-De} & L'utente deve poter impostare il criterio "Parola chiave" per cercare un documento & 
                \ucref{uc:\mdparamParChi} \\
                \hline
                \textbf{R38-F-De} & L'utente deve poter impostare il criterio "Numero allegati" per cercare un documento & 
                \ucref{uc:\mdparamNumAll} \\
                \hline
                \textbf{R39-F-De} & L'utente deve poter impostare il criterio "Impronta crittografica allegato" per cercare un documento & 
                \ucref{uc:\mdparamImpCritAll} \\
                \hline
                \textbf{R40-F-De} & L'utente deve poter impostare il criterio "Identificativo allegato" per cercare un documento & 
                \ucref{uc:\mdparamIdDocAll} \\
                \hline
                \textbf{R41-F-De} & L'utente deve poter impostare il criterio "Descrizione allegato" per cercare un documento & 
                \ucref{uc:\mdparamDesAll} \\
                \hline
                \textbf{R42-F-De} & L'utente deve poter impostare il criterio "Indice di classificazione" per cercare un documento & 
                \ucref{uc:\mdparamIndCla} \\
                \hline
                \textbf{R43-F-De} & L'utente deve poter impostare il criterio "Descrizione dell'indice di classificazione" per cercare un documento & 
                \ucref{uc:\mdparamDesCla} \\
                \hline
                \textbf{R44-F-De} & L'utente deve poter impostare il criterio "Piano di classificazione" per cercare un documento & 
                \ucref{uc:\mdparamPiaCla} \\
                \hline
                \textbf{R45-F-De} & L'utente deve poter impostare il criterio "Riservato" per cercare un documento & 
                \ucref{uc:\mdparamRis} \\
                \hline
                \textbf{R46-F-Ob} & L'utente deve poter impostare il criterio "Formato" per cercare un documento & 
                \ucref{uc:\mdparamFor} \\
                \hline
                \textbf{R47-F-De} & L'utente deve poter impostare il criterio "Nome prodotto software" per cercare un documento & 
                \ucref{uc:\mdparamNomPro} \\
                \hline
                \textbf{R48-F-De} & L'utente deve poter impostare il criterio "Versione prodotto software" per cercare un documento & 
                \ucref{uc:\mdparamVerPro} \\
                \hline
                \textbf{R49-F-De} & L'utente deve poter impostare il criterio "Produttore software" per cercare un documento & 
                \ucref{uc:\mdparamProd} \\
                \hline
                \textbf{R50-F-De} & L'utente deve poter impostare il criterio "Firmato digitalmente" per cercare un documento & 
                \ucref{uc:\mdparamFirDig} \\
                \hline
                \textbf{R51-F-De} & L'utente deve poter impostare il criterio "Sigillato elettronicamente" per cercare un documento & 
                \ucref{uc:\mdparamSigEle} \\
                \hline
                \textbf{R52-F-De} & L'utente deve poter impostare il criterio "Marcatura temporale" per cercare un documento & 
                \ucref{uc:\mdparamMarTem} \\
                \hline
                \textbf{R53-F-De} & L'utente deve poter impostare il criterio "Conformità copie immagine su supporto informatico" per cercare un documento & 
                \ucref{uc:\mdparamConCop} \\
                \hline
                \textbf{R54-F-Ob} & L'utente deve poter impostare il criterio "Tipo aggregazione" per cercare un documento & 
                \ucref{uc:\mdparamTipAgg} \\
                \hline
                \textbf{R55-F-De} & L'utente deve poter impostare il criterio "Identificativo aggregazione" per cercare un documento & 
                \ucref{uc:\mdparamIdAgg} \\
                \hline
                \textbf{R56-F-De} & L'utente deve poter impostare il criterio "Tipologia fascicolo" per cercare un documento & 
                \ucref{uc:\mdparamTipFas} \\
                \hline
                \textbf{R57-F-Ob} & L'utente deve poter impostare il criterio "Nome del documento" per cercare un documento & 
                \ucref{uc:\mdparamNomDoc} \\
                \hline
                \textbf{R58-F-De} & L'utente deve poter impostare il criterio "Versione del documento" per cercare un documento & 
                \ucref{uc:\mdparamVerDoc} \\
                \hline
                \textbf{R59-F-De} & L'utente deve poter impostare il criterio "Tipo modifica" per cercare un documento & 
                \ucref{uc:\mdparamTipMod} \\
                \hline
                \textbf{R60-F-De} & L'utente deve poter impostare il criterio "Data modifica" per cercare un documento & 
                \ucref{uc:\mdparamDatMod} \\
                \hline
                \textbf{R61-F-De} & L'utente deve poter impostare il criterio "Ora modifica" per cercare un documento & 
                \ucref{uc:\mdparamOraMod} \\
                \hline
                \textbf{R62-F-De} & L'utente deve poter impostare il criterio "Tempo di conservazione" per cercare un documento & 
                \ucref{uc:\mdparamTemCon} \\
                \hline
                \textbf{R63-F-De} & L'utente deve poter impostare il criterio "Note" per cercare un documento & 
                \ucref{uc:\mdparamNot} \\
                \hline
                \textbf{R64-F-Ob} & L'utente deve poter cercare un documento impostando uno o più criteri di ricerca & 
                \ucref{uc:\mdsearch} \\
                \hline
                \textbf{R65-F-Ob} & L'utente deve poter visualizzare i riusltati della ricerca, cioè i documenti che soddisfano i criteri di ricerca impostati & 
                \ucref{uc:\searchres} \\
                \hline
                \textbf{R66-F-Ob} & L'utente deve poter visualizzare le informazioni relative ai singoli documenti che rappresentano il risultato della ricerca & 
                \ucref{uc:\searchres} \newline
                \ucref{uc:\searchres.1} \\
                \hline
                \textbf{R67-F-Ob} & L'utente deve ricevere un messaggio informativo se il risultato della ricerca è vuoto, cioè se nessun documento soddisfa i criteri di ricerca impostati & 
                \ucref{uc:\searchres} \newline
                \ucref{uc:\nosearchres} \\
                \hline
                \textbf{R68-F-Ob} & L'utente deve poter visualizzare l'anteprima di un documento selezionato & 
                \ucref{uc:\docpreview} \\
                \hline
                \textbf{R69-F-Ob} & L'utente deve ricevere un messaggio di avviso sull'eventuale indisponibilità dell'anteprima per un documento selezionato & 
                \ucref{uc:\docpreview} \newline
                \ucref{uc:\nodocpreview} \\
                \hline
                \textbf{R70-F-Ob} & L'utente deve poter esportare un documento selezionato, scegliendo il percorso dove salvarlo & 
                \ucref{uc:\exportdoc} \newline
                \ucref{uc:\exportdoc.1} \\
                \hline
                \textbf{R71-F-Ob} & L'utente deve ricevere un messaggio di errore se l'esportazione fallisce & 
                \ucref{uc:\exportdoc} \newline
                \ucref{uc:\exportdocfail} \\
                \hline
                \textbf{R72-F-Op} & L'utente deve poter configurare la ricerca semantica inserendo una query in linguaggio naturale & 
                \ucref{uc:\semparam} \newline
                \ucref{uc:\semparam.1} \\
                \hline
                \textbf{R73-F-Op} & L'utente deve poter visualizzare il risultato della ricerca semantica configurata & 
                \ucref{uc:\semparam} \newline
                \ucref{uc:\semsearch} \\
                \hline
                \textbf{R74-F-Op} & L'utente deve poter stampare un documento selezionato & 
                \ucref{uc:\printdoc} \\
                \hline
                \textbf{R75-F-Ob} & L'utente deve ricevere un messaggio di errore se la stampa fallisce & 
                \ucref{uc:\printdoc} \newline
                \ucref{uc:\printdocfail} \\
                \hline
                \textbf{R76-F-Op} & L'utente deve poter verificare la firma digitale di un documento selezionato & 
                \ucref{uc:\verdigsign} \\
                \hline
                \textbf{R77-F-Op} & L'utente deve poter visualizzare il risultato della verifica della firma digitale di un documento & 
                \ucref{uc:\viewverdigsign} \\
                \hline
            \end{longtable}
        }
        \subsection{Requisiti di Qualità} {
            \rowcolors{2}{lightgray}{white}
            \begin{longtable}{|l|P{0.5\linewidth}|P{0.3\linewidth}|}
                \caption{Requisiti di Qualità} \\
                
                \hline
                \rowcolor{lightgray}
                \textbf{Codice} & \textbf{Descrizione} & \textbf{Fonti} \\
                \hline
                \endfirsthead
                
                \hline
                \rowcolor{lightgray}
                \textbf{Codice} & \textbf{Descrizione} & \textbf{Fonti} \\
                \hline
                \endhead

                \textbf{R1-Q-Ob} & Vengono richiesti i test di unità e di integrazione per il reperimento e l’elaborazione dei dati & \hyperlink{capitolato}{Capitolato di Progetto}, \newline 
                Sez. "Raccomandazioni esplicite" \\
                \hline
                \textbf{R2-Q-Ob} & Sono richiesti i diagrammi UML relativi agli use cases di progetto & \hyperlink{capitolato}{Capitolato di Progetto}, \newline 
                Sez. "Materiale da consegnare" \\
                \hline
                \textbf{R3-Q-Op} & È richiesto lo schema design relativo alla base dati (se utilizzata) o alle API richiamate (se utilizzate) & \hyperlink{capitolato}{Capitolato di Progetto}, \newline 
                Sez. "Materiale da consegnare" \\
                \hline
                \textbf{R4-Q-Ob} & È richiesta la documentazione relativa ai casi di test gestiti e la relativa reportistica & \hyperlink{capitolato}{Capitolato di Progetto}, \newline 
                Sez. "Materiale da consegnare" \\
                \hline
                \textbf{R5-Q-Ob} & È richiesta la lista dei bug risolti durante le fasi di sviluppo & \hyperlink{capitolato}{Capitolato di Progetto}, \newline 
                Sez. "Materiale da consegnare" \\
                \hline
            \end{longtable}
        }
        \subsection{Requisiti di Vincolo} {
            \rowcolors{2}{lightgray}{white}
            \begin{longtable}{|l|P{0.5\linewidth}|P{0.3\linewidth}|}
                \caption{Requisiti di Vincolo} \\
                
                \hline
                \rowcolor{lightgray}
                \textbf{Codice} & \textbf{Descrizione} & \textbf{Fonti} \\
                \hline
                \endfirsthead
                
                \hline
                \rowcolor{lightgray}
                \textbf{Codice} & \textbf{Descrizione} & \textbf{Fonti} \\
                \hline
                \endhead
                
                \textbf{R1-V-Ob} & È richiesto che il sistema funzioni anche in assenza di connettività & \hyperlink{capitolato}{Capitolato di Progetto}, \newline 
                Sez. "Obiettivo" \\
                \hline
                \textbf{R2-V-Ob} & È richiesto che il sistema sia autoconsistente, cioè non deve essere necessario installare nulla sul dispositivo dell’utente & \hyperlink{capitolato}{Capitolato di Progetto}, \newline 
                Sez. "Raccomandazioni esplicite" \\
                \hline
                \textbf{R3-V-Op} & È richiesto l'utilizzo di SQLite come database relazionale & \hyperlink{capitolato}{Capitolato di Progetto}, \newline 
                Sez. "Tecnologie consigliate" \\
                \hline
                \textbf{R4-V-Op} & È richiesto l'utilizzo di Faiss come libreria per la ricerca dei documenti & \hyperlink{capitolato}{Capitolato di Progetto}, \newline 
                Sez. "Tecnologie consigliate" \\
                \hline
                \textbf{R5-V-Ob} & È richiesto l'utilizzo di Angular o React come framework frontend & \hyperlink{capitolato}{Capitolato di Progetto}, \newline 
                Sez. "Tecnologie consigliate" \\
                \hline
                \textbf{R6-V-Ob} & È richiesto l'utilizzo di Github o Bitbucket come sistema di versionamento del codice & \hyperlink{capitolato}{Capitolato di Progetto}, \newline 
                Sez. "Materiale da consegnare" \\
                \hline
            \end{longtable}
        }
    }

\end{document}